\begin{bt}%[1D5H2-3]
% nguon: Archive/2024/TOAN_11_XE/TOAN_11_XE/data4/11C3/KNTT-3.OTC.tex #858
 
	Mẫu số liệu dưới đây ghi lại độ dài quãng đường di chuyển trong một tuần (đơn vị: kilômét) của $40$ chiếc ô tô:
	\begin{center}
		\begin{tabular}{|c|c|c|c|c|c|c|c|c|c|}
			\hline
			100 & 105 & 115 & 116 & 130 & 135 & 138 & 132 & 135 & 120 \\
			\hline
			125 & 128 & 120 & 124 & 140 & 140 & 146 & 145 & 142 & 142 \\
			\hline
			145 & 148 & 150 & 150 & 159 & 155 & 151 & 156 & 155 & 151 \\
			\hline
			154 & 152 & 153 & 160 & 162 & 175 & 176 & 165 & 188 & 198 \\
			\hline
		\end{tabular}
	\end{center}
	\begin{enumerate}[a)]
		\item Lập bảng tần số ghép nhóm bao gồm cả tần số tích lũy với năm nhóm ứng với năm nửa khoảng:
		$$[100;120),[120;140),[140;160),[160;180),[180;200).$$
		\item Xác định số trung bình cộng, trung vị, tứ phân vị của mẫu số liệu ghép nhóm trên.
		\item Mốt của mẫu số liệu ghép nhóm trên là bao nhiêu?
	\end{enumerate}
	\loigiai{
		\begin{center}
			\begin{tabular}{|c|c|c|}
				\hline
				Nhóm & Tần số & Tần số tích lũy \\
				\hline
				$[100;120)$ & $4$ & $4$ \\
				\hline
				$[120;140)$ & $10$ & $14$ \\
				\hline
				$[140;160)$ & $19$ & $33$ \\
				\hline
				[160;180) & $5$ & $38$ \\
				\hline
				$[180;200)$ & $2$ & $40$ \\
				\hline
				\multicolumn{3}{|c|}{$n=40$} \\
				\hline
			\end{tabular}
		\end{center}
		\begin{enumerate} [a)]
			\item Số phần tử của mẫu là $n=40$. Ta có:\\
			$$\dfrac{n}{2}=\dfrac{40}{2}=20$$
			mà $cf_2=14<20<cf_3=33$. Suy ra nhóm $3$ là nhóm đầu tiên có tần số tích lũy lớn hơn hoặc bằng $20$.\\
			Xét nhóm $3$ là nhóm $[140;160)$ có $r=140;d=20;n_3=19$ và nhóm $2$ là nhóm $[140;160)$ có $cf_2=14$.\\
			Áp dụng công thức, ta có trung vị của mẫu số liệu là:
			$$M_e=140+\left(\dfrac{20-14}{19}\right).20 \approx 146$$
			\item Số phần tử của mẫu là $n=40$.
			\begin{itemize} 
				\item Ta có: $\dfrac{n}{4}=10$ mà $4<10 \le 10$. Suy ra nhóm $2$ là nhóm đầu tiên có tần số tích lũy lớn hơn hoặc bằng $10$. Xét nhóm $2$ là nhóm $[120;140)$ có $s=120;h=20;n_2=10$ và nhóm $1$ là nhóm $[100;120)$ có $cf_1=4$.\\
				Áp dụng công thức, ta có tứ phân vị thứ nhất là:
				$$Q_1=120+\left(\dfrac{10-4}{10}\right).20 \approx 132.$$
				\item Ta có tứ phân vị thứ hai là: 
				$$Q_2=M_e \approx 146.$$
				\item Ta có $\dfrac{3n}{4}=30$ mà $14<30<33$. Suy ra nhóm $3$ là nhóm đầu tiên có tần số tích lũy lớn hơn hoặc bằng $30$. Xét nhóm $3$ là nhóm $[140;160)$ có $t=140;l=20;n_3=19$ và nhóm $2$ là nhóm $[120;140)$ có $cf_2=14$.\\
				Áp dụng công thức, ta có tứ phân vị thứ ba là:
				$$Q_3=140+\left(\dfrac{30-14}{19}\right).20 \approx 157.$$	
			\end{itemize}
			\item Ta thấy nhóm $3$ ứng với nửa khoảng $[140;160)$ là nhóm có tần số lớn nhất với $u=140;g=20;n_3=23$. Nhóm $2$ có tần số $n_2=5$, nhóm $4$ có tần số $n_4=6$.\\
			Áp dụng công thức, ta có mốt của mẫu số liệu là:
			$$ M_o=140+\left(\dfrac{23-5}{2.23-5-6}\right) \cdot 20 \approx 170.$$
		\end{enumerate}
	}
\end{bt}

\begin{bt}%[1D5H1-4]
% nguon: Archive/2024/TOAN_11_XE/TOAN_11_XE/data4/11C3/KNTT-3.OTC.tex #863
%[TeX hóa SGK KNTT]%[Phạm Hoài] 
	Người ta ghi lại tuổi thọ của một số con ong cho kết quả như sau
	\begin{center}
		\begin{tabular}{|c|c|c|c|c|c|}
			\hline 
			Tuổi thọ (ngày) & $[0;20)$ &$[20;40)$&$[40;60)$&$[60;80)$&  $[80;100)$\\
			\hline 
			Số lượng&$5$&$12$&$23$&$31$&$29$\\
			\hline 
		\end{tabular}
	\end{center}
	
	Tìm mốt của mẫu số liệu. Giải thích ý nghĩa của giá trị nhận được.
	\loigiai{
		Tần số lớn nhất là $31$ nên nhóm chứa mốt là nhóm $[60;80)$.\\
		Ta có $j=4$; $a_4=60$; $m_4=31$; $m_3=23$; $m_5=29$; $h=20$.\\
		Khi đó $M_O=60+\dfrac{31-23}{(31-23)+(31-29)}\cdot 20=76$.\\
		Như vậy, số con ong có tuổi thọ $76$ ngày là nhiều nhất.	
	} 
\end{bt}

\begin{bt}%[1D5H1-2]
% nguon: Archive/2024/TOAN_11_XE/TOAN_11_XE/data4/11C3/KNTT-3.1.tex #865
%[TeX hóa SGK KNTT]%[Ngọc Hiếu]
	Mẫu số liệu sau cho biết phân bố theo độ tuổi của dân số Việt Nam năm $2019$.
	\begin{center}
		\begin{tabular}{|c|c|c|c|}
			\hline
			Độ tuổi&Dưới $15$ tuổi&Từ $15$ đến dưới $65$ tuổi&Từ $65$ tuổi trở lên\\
			\hline
			Số người& $23\;371\;882$&$65\;420\;451$&$7\;416\;651$\\
			\hline
		\end{tabular}
	\end{center}
	\begin{enumerate}
		\item Mẫu số liệu đã cho có là mẫu số liệu ghép nhóm hay không?
		\item Nêu các nhóm và tần số tương ứng. Dân số Việt Nam năm $2019$ là bao nhiêu?
	\end{enumerate}
	\loigiai{
		\begin{enumerate}
			\item Mẫu số liệu đã cho là mẫu số liệu ghép nhóm.
			\item Có ba nhóm là: Dưới $15$ tuổi, Từ $15$ đến dưới $65$ tuổi, Từ $65$ tuổi trở lên. Có $23\;371\;882$ người dưới $15$ tuổi; $65\;420\;451$ người từ $15$ đến dưới $65$ tuổi và $7\;416\;651$ người từ $65$ tuổi trở lên.\\
			Dân số Việt Nam năm $2019$ là $23\;371\;882+65\;420\;451+7\;416\;651=96\;208\;984$ người. 
		\end{enumerate}
	}
\end{bt}

\begin{bt}%[1D5H1-2]
% nguon: Archive/2024/TOAN_11_XE/TOAN_11_XE/data4/11C3/KNTT-3.1.tex #867

	Mẫu số liệu sau cho biết kết quả kiểm tra môn Toán của lớp $11$A năm 2022.
	\begin{center}
		\begin{tabular}{|c|c|c|c|c|}
			\hline
			Điểm số & $[3;5)$&$[5;7)$&$[7;9)$&$[9;11)$\\
			\hline
			Số học sinh &5&18&10&7\\
			\hline
		\end{tabular}
	\end{center}
	\begin{enumerate}
		\item Mẫu số liệu đã cho có là mẫu số liệu ghép nhóm hay không?
		\item Nêu các nhóm và tần số tương ứng. Số học sinh của lớp 11A là bao nhiêu?
	\end{enumerate}
	\loigiai{
		\begin{enumerate}
			\item Mẫu số liệu đã cho là mẫu số liệu ghép nhóm.
			\item Có bốn nhóm là: Từ 3 đến dưới 5 điểm, từ 5 đến dưới 7 điểm, từ 7 đến dưới 9 điểm, từ 9 đến dưới 11 điểm.\\	
			Có 5 em từ 3 đến dưới 5 điểm; có 18 em từ 5 đến dưới 7 điểm; có 10 em từ 7 đến dưới 9 điểm; có 7 em từ 9 đến dưới 11 điểm.\\
			Số học sinh của lớp 11A là $5+18+10+7=40$ em. 
		\end{enumerate}
	}
\end{bt}

\begin{bt}%[1D5N1-2]
% nguon: Archive/2024/TOAN_11_XE/TOAN_11_XE/data4/11C3/KNTT-3.1.tex #871

	Tính giá trị đại diện và độ dài của mỗi nhóm trong mẫu số liệu sau
	\begin{center}
		\begin{tabular}{|c|c|c|c|c|c|}
			\hline \textbf{Khoảng tuổi} &{[20; 30)} &{[30; 40)} &{[40; 50)} &{[50; 60)} &{[60; 70)} \\
			\hline \textbf{Số khách hàng nữ} & 3 &9 &6 &4 &4 \\
			\hline
		\end{tabular}
	\end{center}
	\loigiai{
		\begin{center}
			\begin{tabular}{|c|c|c|c|c|c|}
				\hline \textbf{Khoảng tuổi} &{[20; 30)} &{[30; 40)} &{[40; 50)} &{[50; 60)} &{[60; 70)} \\
				\hline \textbf{Giá trị đại diện} & 25 & 35 & 45 & 55 & 65 \\
				\hline \textbf{Độ dài của nhóm} & 10 & 10 & 10 & 10 & 10 \\
				\hline
			\end{tabular}
		\end{center}
	}
\end{bt}

\begin{bt}%[1D5H1-2]
% nguon: Archive/2024/TOAN_11_XE/TOAN_11_XE/data4/11C3/KNTT-3.1.tex #872

	Cân nặng của $28$ học sinh nam lớp $11$ được cho như sau:
	\begin{center}
		\begin{tabular}{lllllllllllll}
			$55{,}4$ & $62{,}6$ & $54{,}2$ & $56{,}8$ & $58{,}8$ & $59{,}4$ & $60{,}7$ & $58$ & $59{,}5$ & $63{,}6$ & $61{,}8$ & $52{,}3$ & $63{,}4$ \\
			$49{,}7$ & $45{,}1$ & $56{,}2$ & $63{,}2$ & $46{,}1$ & $49{,}6$ & $59{,}1$ & $55{,}3$ & $55{,}8$ & $45{,}5$ & $46{,}8$ & $54$ & $49{,}2$ \\
			$57{,}9$ & $52{,}6$
		\end{tabular}
	\end{center}
	Hãy chia mẫu dữ liệu trên thành 5 nhóm, lập bảng tần số ghép nhóm và xác định giá trị đại diện cho mỗi nhóm.
	\loigiai{
		Khoảng biến thiên của mẫu số liệu trên là $R=63{,}6-45{,}1=18{,}5$.
		Độ dài mỗi nhóm $L>\dfrac{R}{k}=\dfrac{18{,}5}{5}=3{,}7$.
		Ta chọn $L=4$ và chia dữ liệu thành các nhóm $[45; 49)$, $[49; 53)$, $[53; 57)$, $[57; 61)$, $[61; 65)$.
		Khi đó ta có bảng tần số ghép nhóm sau:
		\begin{center}
			\begin{tabular}{|c|c|c|c|c|c|}
				\hline \textbf{Cân nặng} &{$[45; 49)$} &{$[49; 53)$} &{$[53; 57)$} &{$[57; 61)$} &{$[61; 65)$} \\
				\hline \textbf{Giá trị đại diện} & 47 & 51 & 55 & 59 & 63 \\
				\hline \textbf{Số học sinh} & 4 & 5 & 7 & 7 & 5 \\
				\hline
			\end{tabular}
		\end{center}
	}
\end{bt}

\begin{bt}%[1D5N1-2]
% nguon: Archive/2024/TOAN_11_XE/TOAN_11_XE/data4/11C3/KNTT-3.1.tex #874
%[Tex hóa SGK CD, TVN-223]
	\immini
	{
		\textbf{Bảng bên} biểu diễn mẫu số liệu ghép nhóm được cho dưới dạng bảng tần số ghép nhóm. Hãy cho biết 
		\begin{enumerate}
			\item Mẫu số liệu có bao nhiêu số liệu; bao nhiêu nhóm?
			\item Tần số của mỗi nhóm.
		\end{enumerate}
	}
	{
		\begin{tabular}{|c|c|}
			\hline
			\textbf{Nhóm} & \textbf{Tần số}\\ 
			\hline
			$\left[0;5\right)$ & $11$\\
			\hline
			$\left[5;10\right)$ & $31$\\
			\hline
			$\left[10;15\right)$ & $45$\\
			\hline
			$\left[15;20\right)$ & $21$\\
			\hline
			$\left[20;26\right)$ & $12$\\
			\hline
			& $n = 120$ \\
			\hline
		\end{tabular}
	}
	\loigiai{
		\begin{enumerate}
			\item Mẫu số liệu gồm $120$ số liệu và $5$ nhóm.
			\item Tần số lần lượt của các nhóm $1$, $2$, $3$, $4$, $5$ lần lượt là $11$, $31$, $45$, $21$, $12$.
		\end{enumerate}	
	}
\end{bt}

\begin{bt}%[1D5H1-2]
% nguon: Archive/2024/TOAN_11_XE/TOAN_11_XE/data4/11C3/KNTT-3.1.tex #875
%[Tex hóa SGK CD, TVN-223]
	Một trường trung học phổ thông chọn $36$ học sinh nam của khối $11$, đo chiều cao của các bạn học sinh đó và thu được mẫu số liệu sau (đơn vị: centimét):
	$$
	\begin{array}{llllllllllll}
		160 & 161 & 161 & 162 & 162 & 162 & 163 & 163 & 163 & 164 & 164 & 164 \\
		164 & 165 & 165 & 165 & 165 & 165 & 166 & 166 & 166 & 166 & 167 & 167 \\
		168 & 168 & 168 & 168 & 169 & 169 & 170 & 171 & 171 & 172 & 172 & 174
	\end{array}
	$$
	Lập bảng tần số ghép nhóm cho mẫu số liệu trên có $5$ nhóm ứng với $5$ nửa khoảng:
	$$
	\left[160;163 \right),\ \left[163;169 \right),\ \left[166;169 \right),\ \left[169;172 \right),\ \left[172;175 \right).
	$$
	\loigiai{
		Bảng tần số ghép nhóm như sau:
		\begin{center}
			\begin{tabular}{|c|c|}
				\hline
				\textbf{Nhóm} & \textbf{Tần số}\\ 
				\hline
				$\left[169;163\right)$ & $6$\\
				\hline
				$\left[163;166\right)$ & $12$\\
				\hline
				$\left[166;169\right)$ & $10$\\
				\hline
				$\left[169;172\right)$ & $5$\\
				\hline
				$\left[172;175\right)$ & $3$\\
				\hline
				& $n = 36$ \\
				\hline
			\end{tabular}
		\end{center}
	}
\end{bt}

\begin{bt}%[1D5H1-2]
% nguon: Archive/2024/TOAN_11_XE/TOAN_11_XE/data4/11C3/KNTT-3.1.tex #877

	Trong các mẫu số liệu sau, mẫu nào là mẫu số liệu ghép nhóm? Đọc và giải thích mẫu số liệu ghép nhóm đó.
	\begin{enumerate}
		\item Số tiền lương mà giáo viên nhận được mỗi tháng tại một trường.
		\begin{center}
			\begin{tabular}{|c|c|c|c|c|c|}
				\hline
				Số tiền (triệu đồng)&$[2;4)$&$[4;6)$&$[6;8)$&$[8;10)$&$[10;12)$\\
				\hline
				Số giáo viên&$10$&$50$&$23$&$12$&$5$\\
				\hline
			\end{tabular}
		\end{center}
		\item Thống kê thời gian giải một bài toán khó của 40 học sinh, ta có bảng số liệu sau
		\begin{center}
			\begin{tabular}{|c|c|c|c|c|}
				\hline
				Thời gian (phút)&$[5;9)$&$[9;13)$&$[13;17)$&$[17;21)$\\
				\hline
				Số học sinh &$3$&$12$&$10$&$15$\\
				\hline
			\end{tabular}
		\end{center}
	\end{enumerate}
	\loigiai{
		Cả hai mẫu số liệu trên đều là mẫu số liệu ghép nhóm.
		\begin{enumerate}
			\item Có năm nhóm là
			\begin{enumerate}[\bf --]
				\item Từ $2$ đến dưới $4$ triệu đồng có $10$ giáo viên.
				\item Từ $4$ đến dưới $6$ triệu đồng có $50$ giáo viên.
				\item Từ $6$ đến dưới $8$ triệu đồng có $23$ giáo viên.
				\item Từ $8$ đến dưới $10$ triệu đồng có $12$ giáo viên.
				\item Từ $10$ đến dưới $12$ triệu đồng có $5$ giáo viên.
			\end{enumerate}
			\item Có bốn nhóm là
			\begin{enumerate}[\bf --]
				\item Từ $5$ phút đến dưới $9$ phút có $3$ học sinh.
				\item Từ $9$ phút đến dưới $13$ phút có $12$ học sinh.
				\item Từ $13$ phút đến dưới $17$ phút có $10$ học sinh.
				\item Từ $17$ phút đến dưới $21$ phút có $15$ học sinh.
			\end{enumerate}
		\end{enumerate}
	}
\end{bt}

\begin{bt}%[1D5H1-2]
% nguon: Archive/2024/TOAN_11_XE/TOAN_11_XE/data4/11C3/KNTT-3.1.tex #878

	Cân nặng của $28$ học sinh nam lớp $11$ được cho như sau:
	\begin{center}
		\begin{tabular}{llllllllllllll}
			$55{,}4$ & $62{,}6$ & $54{,}2$ & $56{,}8$ & $58{,}8$ & $59{,}4$ & $60{,}7$ & $58$ & $59{,}5$ & $63{,}6$ & $61{,}8$ & $52{,}3$ & $63{,}4$ & $57{,}9$\\
			$49{,}7$ & $45{,}1$ & $56{,}2$ & $63{,}2$ & $46{,}1$ & $49{,}6$ & $59{,}1$ & $55{,}3$ & $55{,}8$ & $45{,}5$ & $46{,}8$ & $54$ & $49{,}2$ & $52{,}6$
		\end{tabular}
	\end{center}
	Hãy chuyển mẫu số liệu trên sang mẫu số liệu ghép nhóm gồm $5$ nhóm có độ dài bằng nhau.
	\loigiai{
		Khoảng biến thiên của mẫu số liệu trên là $63{,}6-45{,}1=18{,}5$.
		Độ dài mỗi nhóm là $L>\dfrac{18{,}5}{5}=3{,}7$.
		Ta chọn $L=4$ và chia dữ liệu thành các nhóm $[45; 49)$, $[49; 53)$, $[53; 57)$, $[57; 61)$, $[61; 65)$.
		Khi đó ta có bảng tần số ghép nhóm sau:
		\begin{center}
			\begin{tabular}{|c|c|c|c|c|c|}
				\hline Cân nặng &{$[45; 49)$} &{$[49; 53)$} &{$[53; 57)$} &{$[57; 61)$} &{$[61; 65)$} \\
				\hline Số học sinh & 4 & 5 & 7 & 7 & 5 \\
				\hline
			\end{tabular}
		\end{center}
	}
\end{bt}

\begin{bt}%[1D5H1-2]
% nguon: Archive/2024/TOAN_11_XE/TOAN_11_XE/data4/11C3/KNTT-3.1.tex #879

	Anh Văn ghi lại cự li 30 lần ném lao của mình ở bảng sau (đơn vị: mét):
	\begin{center}
		\begin{tabular}{|c|c|c|c|c|c|c|c|c|c|}
			\hline $72{,}1$ & $72{,}9$ & $70{,}2$ & $70{,}9$ & $72{,}2$ & $71{,}5$ & $72{,}5$ & $69{,}3$ & $72{,}3$ & $69{,}7$ \\
			\hline $72{,}3$ & $71{,}5$ & $71{,}2$ & $69{,}8$ & $72{,}3$ & $71{,}1$ & $69{,}5$ & $72{,}2$ & $71{,}9$ & $73{,}1$ \\
			\hline $71{,}6$ & $71{,}3$ & $72{,}2$ & $71{,}8$ & $70{,}8$ & $72{,}2$ & $72{,}2$ & $72{,}9$ & $72{,}7$ & $70{,}7$ \\
			\hline
		\end{tabular}
	\end{center}
	Tổng hợp lại kết quả ném của anh Văn vào bảng tần số ghép nhóm theo mẫu sau:
	\begin{center}
		\begin{tabular}{|c|c|c|c|c|c|}
			\hline Cự li $(\mathrm{m})$ &{$[69{,}2; 70)$} &{$[70; 70{,}8)$} &{$[70{,}8; 71{,}6)$} &{$[71{,}6; 72{,}4)$} &{$[72{,}4; 73{,}2)$} \\
			\hline Số lần & $?$ & $?$ & $?$ & $?$ & $?$ \\
			\hline
		\end{tabular}
	\end{center}
	\loigiai{
		Bảng tần số ghép nhóm kết quả ném của anh Văn:
		\begin{center}
			\begin{tabular}{|c|c|c|c|c|c|}
				\hline Cự li $(\mathrm{m})$ &{$[69{,}2; 70)$} &{$[70; 70{,}8)$} &{$[70{,}8; 71{,}6)$} &{$[71{,}6; 72{,}4)$} &{$[72{,}4; 73{,}2)$} \\
				\hline Số lần & $4$ & $2$ & $7$ & $12$ & $5$ \\
				\hline
			\end{tabular}
		\end{center}
	}
\end{bt}

\begin{bt}%[1D5H1-2]
% nguon: Archive/2024/TOAN_11_XE/TOAN_11_XE/data4/11C3/KNTT-3.1.tex #880

	Một cửa hàng đã thống kê số ba lô bán được mỗi ngày trong tháng 9 với kết quả cho như sau: \begin{center}
		\begin{tabular}{lllllllllllllll}
			$12$ & $29$ & $12$ & $19$ & $15$ & $21$ & $19$ & $29$ & $28$ & $12$ & $15$ & $25$ & $16$ & $20$ & $29$\\
			$21$ & $12$ & $24$ & $14$ & $10$ & $12$ & $10$ & $23$ & $27$ & $28$ & $18$ & $16$ & $10$ & $20$ & $21$
		\end{tabular}
	\end{center}
	Hãy chuyển mẫu số liệu trên sang mẫu số liệu ghép nhóm gồm 5 nhóm có độ dài bằng nhau.
	\loigiai{
		Khoảng biến thiên của mẫu số liệu trên là $29-10=19$.
		Độ dài mỗi nhóm $L>\dfrac{19}{5}=3{,}8$.
		Ta chọn $L=4$ và chia dữ liệu thành các nhóm $[10; 14)$, $[14; 18)$, $[18; 22)$, $[22; 26)$, $[26; 30)$.
		Khi đó ta có bảng tần số ghép nhóm sau:
		\begin{center}
			\begin{tabular}{|c|c|c|c|c|c|}
				\hline Cân nặng &{$[10; 14)$} &{$[14; 18)$} &{$[18; 22)$} &{$[22; 26)$} &{$[26; 30)$} \\
				\hline Số ba lô bán được & $8$ & $5$ & $8$ & $3$ & $6$ \\
				\hline
			\end{tabular}
		\end{center}
	}
\end{bt}

\begin{bt}%[1D5H1-2]
% nguon: Archive/2024/TOAN_11_XE/TOAN_11_XE/data4/11C3/KNTT-3.1.tex #881

	Mẫu số liệu dưới đây ghi lại tốc độ của $40$ ô tô khi đi qua một trạm đo tốc độ (đơn vị: km/h):
	\begin{center}
		\begin{tabular}{llllllllll}
			48{,}5 & 43 & 50 & 55 & 45 & 60 & 53 & 55,5 & 44 & 65 \\
			51 & 62,5 & 41 & 44,5 & 57 & 57 & 68 & 49 & 46{,}5 & 53{,}5 \\
			61 & 49{,}5 & 54 & 62 & 59 & 56 & 47 & 50 & 60 & 61 \\
			49{,}5 & 52{,}5 & 57 & 47 & 60 & 55 & 45 & 47,5 & 48 & 61{,}5
		\end{tabular}
	\end{center}
	Hãy chuyển mẫu số liệu trên sang dạng ghép nhóm với 6 nhóm có độ dài bằng nhau.
	\loigiai{
		Ta có bảng tần số ghép nhóm của mẫu số liệu trên như sau:
		\begin{center}
			\begin{tabular}{|c|c|c|c|c|c|c|}
				\hline Tốc độ (km/h) &
				$\left[40;45\right)$ & 
				$\left[45;50\right)$ & 
				$\left[50;55\right)$ & 
				$\left[55;60\right)$ &
				$\left[60;65\right)$ &
				$\left[65;70\right)$ \\
				\hline Số ôtô & $4$& $11$ & $7$ & $8$ & $8$&$2$\\
				\hline
			\end{tabular}
		\end{center}
	}
\end{bt}

\begin{bt}%[1D5H1-2]
% nguon: Archive/2024/TOAN_11_XE/TOAN_11_XE/data4/11C3/KNTT-3.1.tex #882

	Mẫu số liệu sau ghi lại cân nặng của $30$ bạn học sinh (đơn vị: kilôgam):
	\begin{center}
		\begin{tabular}{llllllllll}
			17 & 40 & 39 & 40{,}5 & 42 & 51 & 41{,}5 & 39 & 41 & 30\\
			40 & 42 & 40{,}5 & 39{,}5 & 41 & 40{,}5 & 37 & 39{,}5 & 40 & 41\\
			38{,}5 & 39{,}5 & 40 & 41 & 39 & 40{,}5 & 40 & 38{,}5 & 39{,}5 & 41{,}5
		\end{tabular}
	\end{center}
	Hãy chuyển mẫu số liệu trên sang dạng ghép nhóm với 8 nhóm có độ dài bằng nhau.
	\loigiai{
		Ta có bảng tần số ghép nhóm của mẫu số liệu trên như sau:
		\begin{center}
			\begin{tabular}{|c|c|c|c|c|c|c|c|c|}
				\hline Tốc độ (km/h) &
				$\left[15;20\right)$ &
				$\left[20;25\right)$ &  
				$\left[25;30\right)$ & 
				$\left[30;35\right)$ & 
				$\left[35;40\right)$ &
				$\left[40;45\right)$ & 
				$\left[45;50\right)$ & 
				$\left[50;55\right)$\\
				\hline Số ôtô & 
				1 &
				0 &
				0 &
				1 &
				10 &
				17 &
				0 &
				1 \\
				\hline
			\end{tabular}
		\end{center}
	}
\end{bt}

\begin{bt}%[1D5H1-3]
% nguon: Archive/2024/TOAN_11_XE/TOAN_11_XE/data4/11C3/KNTT-3.2.tex #884
 
Tìm cân nặng trung bình của học sinh lớp $11D$ cho trong Bảng $3.5$.
\begin{center}
\begin{tabular}{|c|c|c|c|c|c|c|}
	\hline
	Cân nặng	& $\left[40{,}5;45{,}5 \right)$ & $\left[45{,}5;50{,}5 \right)$ & $\left[50{,}5;55{,}5 \right)$ & $\left[55{,}5;60{,}5 \right)$ & $\left[60{,}5;65{,}5 \right)$ & $\left[65{,}5;70{,}5 \right)$ \\
	\hline
Số học sinh&$10$	& $7$ & $16$ &$4$  & $2$ & $3$ \\
	\hline
\end{tabular}

Bảng $3.5$. Cân nặng của học sinh lớp $11D$.	
\end{center}

	\loigiai{
Trong mỗi khoảng cân nặng, giá trị đại diện là trung bình cộng của hai giá trị đầu mút nên ta có bảng sau:
\begin{center}
	\begin{tabular}{|c|c|c|c|c|c|c|}
		\hline
		Cân nặng (kg)	& $43$ & $48$ & $53$ & $58$ & $63$ & $68$ \\
		\hline
		Số học sinh&$10$	& $7$ & $16$ &$4$  & $2$ & $3$ \\
		\hline
	\end{tabular}
\end{center}	
Tổng số học sinh là $n=42$. Cân nặng trung bình của học sinh lớp $11D$ là $$\overline{x}=\dfrac{10\cdot 43+7\cdot 48+16\cdot 53+4\cdot 58+2\cdot 63+3\cdot 68}{42}\approx51{,}81\,\mathrm{(kg)}.$$
	}
\end{bt}

\begin{bt}%[1D5H1-4]
% nguon: Archive/2024/TOAN_11_XE/TOAN_11_XE/data4/11C3/KNTT-3.2.tex #887

Bảng số liệu ghép nhóm sau cho biết chiều cao (cm) của $50$ học sinh lớp $11A$.
\begin{center}
	\begin{tabular}{|c|c|c|c|c|c|c|}
		\hline
		Khoảng chiều cao (cm)	& $\left[145;150 \right)$ & $\left[150;155 \right)$ & $\left[155;160 \right)$ & $\left[160;165 \right)$&$\left[165;170 \right)$  \\
		\hline
		Số học sinh&$7$	& $14$ & $10$ &$10$  & $9$ \\
		\hline
	\end{tabular}
\end{center}
Tính mốt của mẫu số liệu ghép nhóm này. Có thể kết luận gì từ giá trị tính được?	
	\loigiai{
Tần số lớn nhất là $14$ nên nhóm chứa mốt là nhóm $\left[150;155 \right)$. Ta có $j=2$, $a_2=150$, $m_2=14$, $m_1=7$, $m_3=10$, $h=5$. Do đó $$M_o=150+\dfrac{14-7}{\left(14-7\right)+\left(14-10\right)\cdot 5}\approx 153{,}18.$$	
Số học sinh có chiều cao khoảng $153{,}18$ là nhiều nhất.
	}
\end{bt}

\begin{bt}%[1D5H1-3]
% nguon: Archive/2024/TOAN_11_XE/TOAN_11_XE/data4/11C3/KNTT-3.2.tex #888

	Kết quả khảo sát cân nặng của $25$ quả cam ở mỗi lô hàng $A$ và $B$ được cho ở bảng sau:
	\begin{center}
		\begin{tabular}{|c|c|c|c|c|c|}
			\hline \multicolumn{1}{|c|}{Cân nặng $(\mathrm{g})$} &{$[150; 155)$} &{$[155; 160)$} &{$[160; 165)$} &{$[165; 170)$} &{$[170; 175)$} \\
			\hline Số quả cam ở lô hàng $A$ & 2 & 6 & 12 & 4 & 1 \\
			\hline Số quả cam ở lô hàng $B$ & 1 & 3 & 7 & 10 & 4 \\
			\hline
		\end{tabular}
	\end{center}
	\begin{enumerate}
		\item Hãy ước lượng cân nặng trung bình của mỗi quả cam ở lô hàng $A$ và lô hàng $B$.
		\item Nếu so sánh theo số trung bình thì cam ở lô hàng nào nặng hơn?
	\end{enumerate}
	\loigiai{
		Ta có bảng thống kê số lượng cam theo giá trị đại diện:
		\begin{center}
			\begin{tabular}{|c|c|c|c|c|c|}
				\hline \multicolumn{1}{|c|}{Cân nặng $(\mathrm{g})$} &{$152{,}5$} &{$157{,}5$} &{$162{,}5$} &{$167{,}5$} &$172{,}5$\\
				\hline Số quả cam ở lô hàng $A$ & 2 & 6 & 12 & 4 & 1 \\
				\hline Số quả cam ở lô hàng $B$ & 1 & 3 & 7 & 10 & 4 \\
				\hline
			\end{tabular}
		\end{center}
		\begin{enumerate}
			\item Cân nặng trung bình của mỗi quả cam ở lô hàng $A$ xấp xỉ bằng
			\[(2\cdot 152{,}5+6\cdot 157{,}5+12\cdot 162{,}5+4\cdot 167{,}5+1\cdot 172{,}5): 25=161{,}7\ (\mathrm{g}). \]
			Cân nặng trung bình của mỗi quả cam ở lô hàng $B$ xấp xỉ bằng
			\[(1\cdot 152{,}5+3\cdot 157{,}5+7\cdot 162{,}5+10\cdot 167{,}5+4\cdot 172{,}5): 25=165{,}1\ (\mathrm{g}). \]
			\item Nếu so sánh theo số trung bình thì cam ở lô hàng $B$ nặng hơn cam ở lô hàng $A$.
		\end{enumerate}
	}
\end{bt}

\begin{bt}%[1D5V2-2]
% nguon: Archive/2024/TOAN_11_XE/TOAN_11_XE/data4/11C3/KNTT-3.2.tex #890

	Trong tuần lễ bảo vệ môi trường, các học sinh khối $ 11 $ tiến hành thu nhặt vỏ chai nhựa để tái chế. Nhà trường thống kê kết quả thu nhặt vỏ chai của học sinh khối $ 11 $ ở bảng sau:
	\begin{center}
		\begin{tabular}{|c|c|c|c|c|c|}
			\hline 
			\textbf{Số vỏ chai nhựa}	& $ \left[ 11 ; 15\right]  $ & $ \left[ 16 ; 29\right]  $ & $ \left[21 ; 25 \right]  $ & $ \left[ 26 ; 30\right]  $ & $ \left[31 ; 35 \right]  $ \\ 
			\hline 
			\textbf{Số học sinh}	& $ 53 $ & $ 82 $ & $ 48 $ & $ 39 $ & $ 18 $ \\ 
			\hline 
		\end{tabular} 
	\end{center}
	Hãy tìm trung vị của mẫu số liệu ghép nhóm trên.
	\loigiai{
		Do số vỏ chai là số nguyên nên ta hiệu chỉnh lại như sau:
		\begin{center}
			\begin{tabular}{|c|c|c|c|c|c|}
				\hline 
				\textbf{Số vỏ chai nhựa}	& $ \left[ 10{,}5 ; 15{,}5\right) $ & $ \left[ 15{,}5 ; 20{,}5\right) $ & $ \left[ 20{,}5 ; 25{,}5\right) $ & $ \left[ 25{,}5 ; 30{,}5\right) $ & $ \left[30{,}5 ; 35{,}5 \right)  $ \\ 
				\hline 
				\textbf{Số học sinh}& $ 53 $ & $ 82 $ & $ 48 $ & $ 39 $ & $ 18 $ \\ 
				\hline 
			\end{tabular} 
		\end{center}
		Số học sinh tham gia thu nhặt vỏ chai nhựa là $$ n=53+82+48+39+18=240.$$
		Gọi $ x_1; x_2; \ldots ; x_{240} $ lần lượt là số vỏ chai $ 240 $ học sinh khối $ 11 $ thu nhặt được xếp theo thứ tự không giảm.\\
		Do $ x_1, \ldots, x_{53}\in \left[10{,}5 ; 15{,}5 \right) $; $ x_{54}, \ldots, x_{135}\in \left[ 15{,}5 ; 20{,}5\right)$ nên trung vị của mẫu số liệu $ x_1; x_2; \ldots;x_{240} $ là $$ \dfrac{1}{2}\left( x_{120}+x_{121}\right)\in \left[ 15{,}5 ; 20{,}5\right).$$
		Ta xác định được $ n=240$; $ n_m=82 $; $ C=53 $; $ u_m=15{,}5 $; $ u_{m+1}=20{,}5 $.\\
		Trung vị của mẫu số liệu ghép nhóm là $$ M_e=15{,}5+\dfrac{\dfrac{240}{2}-53}{82}\cdot \left( 20{,}5-15{,}5\right)=\dfrac{803}{41}\approx 19{,}59. $$
	}
\end{bt}

\begin{bt}%[1D5V1-2]
% nguon: Archive/2024/TOAN_11_XE/TOAN_11_XE/data4/11C3/KNTT-3.2.tex #891

	Trong một hội thao, thời gian chạy $ 200 $ m của một nhóm các vận động viên được ghi lại ở bảng sau:
	\begin{center}
		\begin{tabular}{|c|c|c|c|c|c|}
			\hline 
			\textbf{Thời gian} \textbf{(giây)}& $ \left[21 ; 21{,}5 \right)  $ & $ \left[ 21{,}5 ; 22\right)  $ & $ \left[ 22 ; 22{,}5\right)  $ & $ \left[ 22{,}5 ; 23\right)  $ & $ \left[ 23 ; 23{,}5\right)  $ \\ 
			\hline 
			\textbf{Số vận động viên} & $ 5 $ & $ 12 $ & $ 32 $ & $ 45 $ & $ 30 $ \\ 
			\hline 
		\end{tabular} 
	\end{center}
	Dựa vào bảng số liệu trên, ban tổ chức muốn chọn ra khoảng $ 50 \% $ số vận động viên chạy nhanh nhất để tiếp tục thi vòng $ 2 $. Ban tổ chức nên chọn các vận động viên có thời gian chạy không quá bao nhiêu giây?
	\loigiai{
		Số vận động viên tham gia là $$n=5+12+32+45+30=124.$$
		Gọi $ x_1; x_2; \ldots ; x_{124} $ lần lượt là thời gian chạy $ 200 $ m của $ 124 $ vận động viên được xếp theo thứ tự không giảm.\\
		Do $ x_1, \ldots, x_5 \in \left[ 21 ; 21{,}5 \right)$, $ x_6, \ldots, x_{17} \in \left[ 21{,}5 ; 22\right) $, $ x_{18}, \ldots, x_{49} \in \left[22 ; 22{,}5\right) $, $ x_{50},\ldots, x_{94} \in \left[ 22{,}5 ; 23\right) $ nên trung vị của mẫu số liệu $ x_1; x_2; \ldots ;x_{124} $ là
		$$\dfrac{1}{2}\cdot \left( x_{62}+x_{63}\right) \in  \left[ 22{,}5 ; 23\right).$$
		Ta xác định được $ n=124 $; $ n_m=45$; $ C=5+12+32=49 $; $ u_m= 22{,}5$; $ u_{m+1}=23$.\\
		Trung vị của mẫu số liệu ghép nhóm là $$M_e=22{,}5 +\dfrac{\dfrac{124}{2}-49}{45}\cdot \left( 23-22{,}5\right)= \dfrac{1019}{45}\approx 22{,}64.$$
		Vậy ban tổ chức nên chọn các vận động viên  có thời gian chạy không quá $ 22{,}64$ (giây) được tiếp tục thi vòng hai.
	}
\end{bt}

\begin{bt}%[1D5H2-3]
% nguon: Archive/2024/TOAN_11_XE/TOAN_11_XE/data4/11C3/KNTT-3.2.tex #892
%1
	Tìm tứ phân vị thứ nhất $Q_1$ và tứ phân vị thứ ba $Q_3$ của mẫu số liệu ghép nhóm cho trong bảng sau
	\begin{center}
		\begin{tabular}{|l|c|c|c|c|c|}
			\hline Thời gian (phút)  &{$[9,5 ; 12,5)$}&{$[12,5 ; 15,5)$}&{$[15,5 ; 18,5)$}&{$[18,5 ; 21,5)$}&{$[21,5 ; 24,5)$}\\
			\hline Số học sinh  & $3$ & $12$ & $15$ & $24$ & $2$ \\
			\hline
		\end{tabular}     
	\end{center}
	\loigiai{
		Cỡ mẫu là $n=56$.\\
		Tứ phân vị thứ nhất $Q_1$ là $\dfrac{x_{14}+x_{15}}{2}$. Do $x_{14}$, $x_{15}$ đều thuộc nhóm $[12,5 ; 15,5)$ nên nhóm này chứa $Q_1$. \\Do đó, $p=2 ; \;a_2=12,5 ;\; m_2=12 ; \;m_1=3, \;a_3-a_2=3$ và ta có
		$$
		Q_1=12,5+\dfrac{\frac{56}{4}-3}{12}\cdot 3=15,25.
		$$
		Với tứ phân vị thứ ba $Q_3$ là $\dfrac{x_{42}+x_{43}}{2}$. Do $x_{42},\; x_{43}$ đều thuộc nhóm $[18,5 ; 21,5)$ nên nhóm này chứa $Q_3$. Do đó, $p=4 ; \;a_4=18,5 ;\; m_4=24 ; \;m_1+m_2+m_3=3+12+15=30 ; \;a_5-a_4=3$ và ta có
		$$
		Q_3=18,5+\dfrac{\frac{3 \cdot 56}{4}-30}{24}\cdot 3=20.
		$$
	}
\end{bt}

\begin{bt}%[1D5H2-3]
% nguon: Archive/2024/TOAN_11_XE/TOAN_11_XE/data4/11C3/KNTT-3.2.tex #893
%2
	Trong phòng thí nghiệm, người ta chia $99$ mẫu vật thành năm nhóm căn cứ trên khối lượng của chúng (đơn vị: gam) và lập bảng tần số ghép nhóm như sau.
	\begin{center}
		\begin{tabular}{|l|c|c|c|c|c|}
			\hline Nhóm (gam)  &{$[27,5 ; 32,5)$}&{$[32,5 ; 37,5)$}&{$[37,5 ; 42,5)$}&{$[42,5 ; 47,5)$}&{$[47,5 ; 52,5)$}\\
			\hline Số mẫu vật  & $16$ & $24$ & $20$ & $30$ & $9$ \\
			\hline
		\end{tabular}     
	\end{center}
	Tìm tứ phân vị của mẫu số liệu.
	\loigiai{
		Cỡ mẫu là $n=99$.\\
		Gọi $x_1, \ldots, x_{99}$ là mẫu số liệu được sắp xếp theo thứ tự không giảm. Khi đó, trung vị là $x_{50}$. Do $x_{50}$ thuộc nhóm $[37,5 ; 42,5)$ nên nhóm này chứa $Q_2$. \\Do đó, $p=3 ; \;a_3=37,5 \;m_3=20;\; m_1+m_2=40 ;  \;a_4-a_3=5$ và ta có
		$$
		Q_2=37,5+\dfrac{\frac{99}{2}-40}{20}\cdot 5=39,875.
		$$
		Tứ phân vị thứ nhất $Q_1$ là $x_{25}$. Do $x_{25}$ đều thuộc nhóm $[32,5 ; 37,5)$ nên nhóm này chứa $Q_1$. \\Do đó, $p=2 ; \;a_2=32,5 ;\; m_2=24 ; \;m_1=16, \;a_3-a_2=5$ và ta có
		$$
		Q_1=32,5+\dfrac{\frac{99}{4}-16}{24}\cdot 5\approx34,323.
		$$
		Với tứ phân vị thứ ba $Q_3$ là $x_{75}$. Do $x_{75}$ đều thuộc nhóm $[42,5 ; 47,5)$ nên nhóm này chứa $Q_3$. Do đó, $p=4 ; \;a_4=42,5 ;\; m_4=30 ; \;m_1+m_2+m_3=3+12+15=60 ; \;a_5-a_4=5$ và ta có
		$$
		Q_3=42,5+\dfrac{\frac{3 \cdot 99}{4}-60}{30}\cdot 5=44,875.
		$$
	}
\end{bt}

\begin{bt}%[1D5H2-3]
% nguon: Archive/2024/TOAN_11_XE/TOAN_11_XE/data4/11C3/KNTT-3.2.tex #894
%3
	Sau khi điều tra về số học sinh trong $100$ lớp học, người ta chia mẫu số liệu đó thành năm nhóm căn cứ vào số lượng học sinh của mỗi lớp (đơn vị: học sinh) và lập bảng lần số ghép nhóm như sau. Tìm tứ phân vị của mẫu số liệu đó.
	\begin{center}
		\begin{tabular}{|l|c|c|c|c|c|}
			\hline Nhóm (học sinh)  &{$[36 ; 38)$}&{$[38 ; 40)$}&{$[40 ; 42)$}&{$[42 ; 44)$}&{$[44 ; 46)$}\\
			\hline Số lớp & $9$ & $15$ & $25$ & $30$ & $21$ \\
			\hline
		\end{tabular}     
	\end{center}
	\loigiai{
		Cỡ mẫu là $n=100$.\\
		Gọi $x_1, \ldots, x_{100}$ là mẫu số liệu được sắp xếp theo thứ tự không giảm. Khi đó, trung vị là $\dfrac{x_{50}+x_{51}}{2}$. Do $x_{50}$, $x_{51}$ thuộc nhóm $[42 ; 44)$ nên nhóm này chứa $Q_2$. \\Do đó, $p=4 ; \;a_4=42 \;m_4=30;\; m_1+m_2+m_3=49 ;  \;a_5-a_4=2$ và ta có
		$$
		Q_2=42+\dfrac{\frac{100}{2}-49}{30}\cdot 2\approx42,06.
		$$
		Tứ phân vị thứ nhất $Q_1$ là $\dfrac{x_{25}+x_{26}}{2}$. Do $x_{25}$, $x_{26}$ đều thuộc nhóm $[40 ; 42)$ nên nhóm này chứa $Q_1$. \\Do đó, $p=3 ; \;a_3=40 ;\; m_3=25 ; \;m_1+m_2=24, \;a_4-a_3=2$ và ta có
		$$
		Q_1=40+\dfrac{\frac{100}{4}-24}{25}\cdot 2\approx40,08.
		$$
		Với tứ phân vị thứ ba $Q_3$ là $\dfrac{x_{75}+x_{76}}{2}$. Do $x_{75}$, $x_{76}$ đều thuộc nhóm $[42 ; 44)$ nên nhóm này chứa $Q_3$. Do đó, $p=4 ; \;a_4=42 ;\; m_4=30 ; \;m_1+m_2+m_3=3+12+15=49 ; \;a_5-a_4=2$ và ta có
		$$
		Q_3=42+\dfrac{\frac{3 \cdot 100}{4}-49}{30}\cdot 2\approx43,7.
		$$
	}
\end{bt}

\begin{bt}%[1D5H2-3]
% nguon: Archive/2024/TOAN_11_XE/TOAN_11_XE/data4/11C3/KNTT-3.2.tex #895
%4
	Giáo viên chủ nhiệm chia thời gian sử dụng Internet trong một ngày của 40 học sinh thành năm nhóm (đơn vị: phút) và lập bảng tẫn số ghép nhóm như sau.
	Tìm tứ phân vị của mẫu số liệu đó.
	\begin{center}
		\begin{tabular}{|l|c|c|c|c|c|}
			\hline Nhóm (phút)  &{$[0 ; 60)$}&{$[60 ; 120)$}&{$[120 ; 180)$}&{$[180 ; 240)$}&{$[240 ; 300)$}\\
			\hline Số học sinh & $6$ & $13$ & $13$ & $6$ & $2$ \\
			\hline
		\end{tabular}     
	\end{center}
	\loigiai{
		Cỡ mẫu là $n=40$.\\
		Gọi $x_1, \ldots, x_{40}$ là mẫu số liệu được sắp xếp theo thứ tự không giảm. Khi đó, trung vị là $\dfrac{x_{20}+x_{21}}{2}$. Do $x_{20}$, $x_{21}$ thuộc nhóm $[120 ; 180)$ nên nhóm này chứa $Q_2$. \\Do đó, $p=3 ; \;a_3=120 \;m_3=13;\; m_1+m_2=19 ;  \;a_4-a_3=60$ và ta có
		$$
		Q_2=120+\dfrac{\frac{40}{2}-19}{13}\cdot 60\approx124,62.
		$$
		Tứ phân vị thứ nhất $Q_1$ là $\dfrac{x_{10}+x_{11}}{2}$. Do $x_{10}$, $x_{11}$ đều thuộc nhóm $[60 ; 120)$ nên nhóm này chứa $Q_1$. \\Do đó, $p=2 ; \;a_2=60 ;\; m_2=13 ; \;m_1=6, \;a_3-a_2=60$ và ta có
		$$
		Q_1=60+\dfrac{\frac{40}{4}-6}{13}\cdot 60\approx78,46.
		$$
		Với tứ phân vị thứ ba $Q_3$ là $\dfrac{x_{30}+x_{31}}{2}$. Do $x_{30}$, $x_{31}$ đều thuộc nhóm $[120 ; 180)$ nên nhóm này chứa $Q_3$. Do đó, $p=3 ; \;a_3=120 ;\; m_3=13 ; \;m_1+m_2=19 ; \;a_4-a_3=60$ và ta có
		$$
		Q_3=120+\dfrac{\frac{3 \cdot 40}{4}-19}{13}\cdot 60\approx170,77.
		$$
	}
\end{bt}

\begin{bt}%[1D5H2-3]
% nguon: Archive/2024/TOAN_11_XE/TOAN_11_XE/data4/11C3/KNTT-3.2.tex #896

	Thời gian luyện tập trong một ngày (tính theo giờ) của một số vận động viên được ghi lại ở bảng sau
	\begin{center}
		\begin{tabular}{|l|c|c|c|c|c|}
			\hline Thời gian luyện tập (giờ)  &{$[0 ; 2)$}&{$[2 ;4)$}&{$[4 ; 6)$}&{$[6 ; 8)$}&{$[8 ; 10)$}\\
			\hline Số vận động viên & $3$ & $8$ & $12$ & $12$ & $4$ \\
			\hline
		\end{tabular}     
	\end{center}
	Tìm tứ phân vị của mẫu số liệu đó.
	\loigiai{
		Cỡ mẫu là $n=39$.\\
		Gọi $x_1, \ldots, x_{39}$ là mẫu số liệu được sắp xếp theo thứ tự không giảm. Khi đó, trung vị là $x_{20}$. Do $x_{20}$ thuộc nhóm $[4 ; 6)$ nên nhóm này chứa $Q_2$. \\Do đó, $p=3 ; \;a_3=4 \;m_3=12;\; m_1+m_2=11 ;  \;a_4-a_3=2$ và ta có
		$$
		Q_2=4+\dfrac{\frac{39}{2}-11}{12}\cdot 2\approx5,417.
		$$
		Tứ phân vị thứ nhất $Q_1$ là $x_{10}$. Do $x_{10}$  thuộc nhóm $[2;4)$ nên nhóm này chứa $Q_1$. \\Do đó, $p=2 ; \;a_2=2 ;\; m_2=8 ; \;m_1=3, \;a_3-a_2=2$ và ta có
		$$
		Q_1=2+\dfrac{\frac{39}{4}-3}{8}\cdot 2=3,6875.
		$$
		Với tứ phân vị thứ ba $Q_3$ là $x_{30}$. Do $x_{30}$  thuộc nhóm $[6;8)$ nên nhóm này chứa $Q_3$. Do đó, $p=3 ; \;a_3=4 ;\; m_3=12 ; \;m_1+m_2=11 ; \;a_4-a_3=2$ và ta có
		$$
		Q_3=4+\dfrac{\frac{3 \cdot 39}{4}-11}{12}\cdot 2\approx7,042.
		$$ 
	}
\end{bt}

\begin{bt}%[1D5H1-4]
% nguon: Archive/2024/TOAN_11_XE/TOAN_11_XE/data4/11C3/KNTT-3.2.tex #898
%2
	Tuồi thọ (năm) của 50 bình ắc quy ô tô được cho như sau:
	\begin{center}
		\begin{tabular}{|l|c|c|c|c|c|c|}
			\hline Tuồi thọ (năm) &{$[2 ; 2,5)$}&{$[2,5 ; 3)$}&{$[3 ; 3,5)$}&{$[3,5 ; 4)$}&{$[4 ; 4,5)$}&{$[4,5 ; 5)$}\\
			\hline Tằn số & $4$ & $9$ & $14$ & $11$ & $7$ & $5$ \\
			\hline
		\end{tabular}
	\end{center}
	Xác định mốt và giải thích ý nghĩa.
	\loigiai{
		Tần số lớn nhất là $14$ nên nhóm chứa mốt là nhóm $[3 ; 3,5)$. \\
		Ta có $j=3, a_3=3, m_3=14$, $m_2=9, m_4=11, h=0,5$. Do đó
		$$
		M_0=3+\frac{14-9}{(14-9)+(14-11)}\cdot 0,5 = 3,3125.
		$$
		Bình ắc quy ô tô có tuổi thọ $3,3125$ năm là nhiều nhất. 
	}
\end{bt}

\begin{bt}%[1D5H1-4]
% nguon: Archive/2024/TOAN_11_XE/TOAN_11_XE/data4/11C3/KNTT-3.2.tex #899
%3
	Một thư viện thống kê số lượng sách được mượn mỗi ngày trong ba tháng ở bảng sau
	\begin{center}
		\begin{tabular}{|l|c|c|c|c|c|c|c|}
			\hline Số sách &{$[16;20]$}&{$[21;25]$}&{$[26;30]$}&{$[31;35]$}&{$[36;40]$}&{$[41;45]$}\\
			\hline Số ngày &$3$ & $6$ & $15$ & $27$ & $22$ & $14$ \\
			\hline
		\end{tabular}
	\end{center}
	Xác định mốt của mẫu số liệu ghép nhóm này.
	\loigiai{Ta hiệu chỉnh bảng số liệu như sau
		\begin{center}
			\begin{tabular}{|l|c|c|c|c|c|c|c|}
				\hline Số sách &{$[15,5;20,5)$}&{$[20,5;25,5)$}&{$[25,5;30,5)$}&{$[30,5;35,5)$}&{$[35,5;40,5)$}&{$[40,5;45,5)$}\\
				\hline Số ngày &$3$ & $6$ & $15$ & $27$ & $22$ & $14$ \\
				\hline
			\end{tabular}
		\end{center} 
		Tần số lớn nhất là $27$ nên nhóm chứa mốt là nhóm $[30,5;35,5)$. \\
		Ta có $j=4, a_4=31, m_4=27$, $m_3=15, m_5=22, h=5$. Do đó
		$$
		M_0=30,5+\frac{27-15}{(27-15)+(27-22)}\cdot 5 \approx 34,03.
		$$
	}
\end{bt}

\begin{bt}%[1D5H1-4]
% nguon: Archive/2024/TOAN_11_XE/TOAN_11_XE/data4/11C3/KNTT-3.2.tex #901

	Điểm kiểm tra môn Toán của lớp 12A được cho trong bảng sau
	\begin{center}
		\begin{tabular}{|l|c|c|c|c|c|c|}
			\hline Thời gian (phút) &{$[3;5)$}&{$[5;7)$}&{$[7;9)$}&{$[9;11)$}\\
			\hline Số học sinh & $5$ & $18$ & $10$ & $7$  \\
			\hline
		\end{tabular}
	\end{center}
	Xác định mốt của mẫu số liệu ghép nhóm này.
	\loigiai{
		Tần số lớn nhất là $18$ nên nhóm chứa mốt là nhóm $[5;7)$. \\
		Ta có $j=2; a_2=5; m_2=18$, $m_3=10, m_1=5, h=2$. Do đó
		$$
		M_0=5+\frac{18-5}{(18-5)+(18-10)}\cdot 2 \approx 6,24.
		$$
	}
\end{bt}

\begin{bt}%[1D5H1-4]
% nguon: Archive/2024/TOAN_11_XE/TOAN_11_XE/data4/11C3/KNTT-3.2.tex #902

Quãng đường (km) từ nhà đến nơi làm việc của 40 công nhân một nhà máy được ghi lại như sau:
\begin{center}
\begin{tabular}{cccccccccccccccccc}
	
$5$	& $3$ &$10$ & $20$ & $25$ & $11$ & $13$ & $7$ & $12$ & $31$ & $19$ &$10$  &$12$  & $17$ & $18$ & $11$ & $32$ & $17$  \\
$7$	& $9$ &$7$ & $8$ & $3$ & $5$ & $12$ & $15$ & $18$ & $3$ & $12$ &$14$  &$2$  & $9$ & $6$ & $15$ & $15$ & $7$  \\
$16$  &$2$ &$6$  &$12$
	
\end{tabular}
\end{center}
\begin{enumerate}
	\item [a)] Ghép nhóm dãy số liệu trên thành các khoảng có độ rộng bằng nhau, khoảng đầu tiên là $\left[0;5\right)$. Tìm giá trị đại diện cho mỗi nhóm.
	\item [b)] Tính số trung bình của mẫu số liệu không ghép nhóm và mẫu số liệu ghép nhóm. Giá trị nào chính xác hơn?
	\item [c)] Xác định nhóm chứa mốt của mẫu số liệu ghép nhóm thu được.
\end{enumerate}
	\loigiai{
\begin{enumerate}
	\item [a)] Giá trị nhỏ nhất của mẫu số liệu là 2, giá trị lớn nhất là 32, khoảng đầu tiên của mẫu số liệu ghép nhóm là $\left[0;5\right)$ nên ta ghép nhóm mẫu số liệu như sau:
	\begin{center}
		\begin{tabular}{|c|c|c|c|c|c|c|c|}
			\hline
	Quãng đường		 & $\left[0;5\right)$ & $\left[5;10\right)$ & $\left[10;15\right)$ & $\left[15;20\right)$ & $\left[20;25\right)$& $\left[25;30\right)$& $\left[30;35\right)$\\
			\hline
	Số công nhân		& $5$ & $11$ & $11$ & $9$ & $1$ & $1$ & $2$ \\
			\hline
		\end{tabular}
	\end{center}
	
	Trong mỗi khoảng, giá trị đại điện là trung bình cộng của hai giá trị đầu mút nên ta có bảng sau:
		\begin{center}
		\begin{tabular}{|c|c|c|c|c|c|c|c|}
			\hline
			Quãng đường		 & $2{,}5$ & $7{,}5$ & $12{,}5$ & $17{,}5$ & $22{,}5$& $27{,}5$& $32{,}5$\\
			\hline
			Số công nhân		& $5$ & $11$ & $11$ & $9$ & $1$ & $1$ & $2$ \\
			\hline
		\end{tabular}
	\end{center}
	
	\item [b)] Số trung bình của mẫu số liệu không ghép nhóm là
	$$\overline{x}=\dfrac{5+3+10+\cdots +12}{40}=11{,}9.$$
	Số trung bình của mẫu số liệu ghép nhóm là
	$$\overline{x}=\dfrac{5\cdot 2{,}5+11\cdot 7{,}5+11\cdot 12{,}5+9\cdot 17{,}5+1\cdot 22{,}5+1\cdot 27{,}5+2\cdot 32{,}5}{40}=12{,}625.$$
	Số trung bình của mẫu số liệu không ghép nhóm sẽ chính xác hơn số trung bình của mẫu số liệu ghép nhóm vì số trung bình của dữ liệu không ghép nhóm sử dụng chính xác các số liệu, còn số trung bình của dữ liệu ghép nhóm sử dụng giá trị đại diện của mỗi khoảng ghép nhóm.
	\item [c)] Tần số lớn nhất là 11 nên nhóm chứa mốt của mẫu số liệu ghép nhóm là nhóm $\left[5;10\right)$ và nhóm $\left[10;15\right)$.
\end{enumerate}
	}
\end{bt}

\begin{bt}%[1D5H1-4]
% nguon: Archive/2024/TOAN_11_XE/TOAN_11_XE/data4/11C3/KNTT-3.2.tex #903

Tuổi thọ (năm) của 50 bình ắc quy ô tô được cho như sau:
\begin{center}
	\begin{tabular}{|c|c|c|c|c|c|c|}
		\hline
		Tuổi thọ (năm)	& $\left[2;2{,}5 \right)$ & $\left[2{,}5;3 \right)$ & $\left[3;3{,}5 \right)$&$\left[3{,}5;4 \right)$&$\left[4;4{,}5 \right)$&$\left[4{,}5;5 \right)$  \\
		\hline
	Tần số &$4$	& $9$ & $14$ &$11$  & $7$&$5$ \\
		\hline
	\end{tabular}
\end{center}
\begin{enumerate}
	\item [a)] Xác định mốt và giải thích ý nghĩa.
	\item [b)] Tính tuổi thọ trung bình của $50$ bình ắc quy ô tô này.
\end{enumerate}	
	\loigiai{
\begin{enumerate}
	\item [a)] Tần số lớn nhất là $14$ nên nhóm chứa mốt là nhóm $\left[3;3, 5\right)$.\\
	Ta có $j=3$, $a_3=3$, $m_3=14$, $m_2=9$, $m_4=11$, $h=0$, $5$. Do đó
	$$M_o=3+\dfrac{14-9}{(14-9)+(14-11)}\cdot 0{,}5=3{,}3125.$$
	Tuổi thọ cao nhất của bình ắc quy khoảng $3{,}3125$ năm.
	\item [b)] Ta có bảng sau
	\begin{center}
	\begin{tabular}{|c|c|c|c|c|c|c|}
	\hline
	Tuổi thọ (năm)	& $2{,}25$ & $2{,}75$ & $3{,}25$&$3{,}75$&$4{,}25$&$4{,}75$  \\
	\hline
	Tần số &$4$	& $9$ & $14$ &$11$  & $7$&$5$\\
	\hline
\end{tabular}	
\end{center}
	Tuổi thọ trung bình của 50 bình ắc quy ô tô này là
	$$\overline{x}=\dfrac{2{,}25\cdot 4+2{,}75\cdot 9+3{,}25\cdot 14+3{,}75\cdot 11+4{,}25\cdot 7+4{,}75\cdot 5}{50}=3{,}48 \, \text{(năm)}.$$
\end{enumerate}	}
\end{bt}

\begin{bt}%[1D5H1-3]
% nguon: Archive/2024/TOAN_11_XE/TOAN_11_XE/data4/11C3/KNTT-3.2.tex #905

\immini{
Phỏng vấn một số học sinh lớp $11$ về thời gian (giờ) ngủ của một buổi tối, thu được bảng số liệu ở bên.
\begin{enumerate}
	\item [a)] So sánh thời gian ngủ trung bình của các bạn học sinh nam và nữ.
	\item [b)] Hãy cho biết $75\%$ học sinh khối $11$ ngủ ít nhất bao nhiêu giờ?
\end{enumerate}
}
{
\begin{tabular}{|c|c|c|}
	\hline
Thời gian	& Số học sinh nam & Số học sinh nữ\\
	\hline
$\left[4;5 \right)$	& $6$ & $4$ \\
	\hline
$\left[5;6 \right)$	& $10$ & $8$ \\
	\hline
$\left[6;7 \right)$	& $13$ & $10$ \\
	\hline
$\left[7;8 \right)$	& $9$ & $11$ \\
	\hline
$\left[8;9 \right)$	& $7$ & $8$ \\
	\hline
\end{tabular}
}
	\loigiai{
	\begin{enumerate}
		\item [a)] Trong mỗi khoảng thời gian, giá trị đại diện là trung bình cộng của giá trị hai đầu mút nên ta có bảng sau:
		\begin{center}
	\begin{tabular}{|c|c|c|}
		\hline
		Thời gian	& Số học sinh nam & Số học sinh nữ\\
		\hline
		$4{,}5$	& $6$ & $4$ \\
		\hline
		$5{,}5$	& $10$ & $8$ \\
		\hline
		$6{,}5$	& $13$ & $10$ \\
		\hline
		$7{,}5$	& $9$ & $11$ \\
		\hline
		$8{,}5$	& $7$ & $8$ \\
		\hline
	\end{tabular}	
	\end{center}
		Tổng số học sinh nam là $n_1=6+10+13+9+7=45$.\\ Thời gian ngủ trung bình của học sinh nam là:
		$$\overline{x_1}=\dfrac{4{,}5\cdot 6+5{,}5\cdot10+6{,}5\cdot13+7{,}5\cdot9+8{,}5\cdot7}{45}=\dfrac{587}{90}\approx 6{,}52\,\, \text{(giờ)}.$$
		Tổng số học sinh nữ là $n_2=4+8+10+11+8=41$. Thời gian ngủ trung bình của học sinh nữ là:
		$$\overline{x_2}=\dfrac{4,5\cdot4+5,5\cdot8+6,5\cdot10+7,5\cdot11+8,5\cdot8}{41}=\dfrac{555}{82}\approx 6{,}77 \,\,\text{(giờ)}.$$
		Vì $\overline{x_2}>\overline{x_1}$ nên thời gian ngủ trung bình của các bạn học sinh nữ lớn hơn thời gian ngủ trung bình của các bạn nam.
		\item [b)] Tổng số học sinh được điều tra là $n=n_1+n_2=45+41=86$.\\
		Giả sử $x_1;x_2;x_3;\cdot \cdot;x_{86}$ là dãy giá trị được sắp xếp theo thứ tự không giảm.\\
		Ta có bảng sau:
		\begin{center}
		\begin{tabular}{|c|c|c|}
			\hline
			Thời gian	& Số học sinh \\
			\hline
			$\left[4;5 \right)$	& $10$  \\
			\hline
			$\left[5;6 \right)$	& $18$  \\
			\hline
			$\left[6;7 \right)$	& $23$  \\
			\hline
			$\left[7;8 \right)$	& $20$  \\
			\hline
			$\left[8;9 \right)$	& $15$  \\
			\hline
		\end{tabular}
	\end{center}
		Tứ phân vị thứ nhất $Q_1$ là $x_{22}$. Do $x_{22}$ thuộc nhóm $\left[5;6\right)$ nên nhóm này chứa $Q_1$.\\ Do đó, $p=2;\,a_2=5;\,m_2=18;\,m_1=10;\,a_3-a_2=1$ và ta có
		$$Q_1=5+\dfrac{\dfrac{86}{4}-10}{18}\cdot 1=\dfrac{203}{36}\approx 5{,}64 \text{(giờ)}.$$
		Nghĩa là có $25\%$ học sinh khối $11$ ngủ ít hơn $5{,}64$ giờ.\\ Vậy $75\%$ học sinh khối $11$ ngủ ít nhất $5{,}64$ giờ.
	\end{enumerate}
	}
\end{bt}

\begin{bt}%[1D5H1-4]
% nguon: Archive/2024/TOAN_11_XE/TOAN_11_XE/data4/11C3/KNTT-3.2.tex #910

	Kết quả đo chiều cao của $200$ cây keo $3$ năm tuổi ở một nông trường được biểu diễn ở biểu đồ dưới đây.
	\begin{center}
		\begin{tikzpicture}[scale=1,font=\scriptsize]
			\def\hoanh{11.5};
			\def\tung{6.5};
			\def\mau{cyan};
			\foreach \x/\n in{1/20,3/35,5/60,7/55,9/30}{\draw[line width=16mm,\mau] (\x,0)--++(0,{\n/10});
				\draw[dashed] (\x,{\n/10})node[above]{$\n$}--(0,{\n/10}) node[left]{$\n$};}
			\foreach \x/\p in {1/[8{,}5;8{,}8),3/[8{,}8;9{,}1),5/[9{,}1;9{,}4),7/[9{,}4;9{,}7),9/[9{,}7;10{,}0)}{\node[below] at (\x,0){\scriptsize $\p$};}
			\draw[->] (0,0)--(\hoanh,0) node[below]{($m$)};
			\draw[->] (0,0)node[below left]{$O$}--(0,\tung) node[left]{(Số cây)};
			\path (current bounding box.north) node[above]		{\textbf{Chiều cao 200 cây keo 3 năm tuổi}};
		\end{tikzpicture}
	\end{center}
	Hãy ước lượng số trung bình và mốt của mẫu số liệu ghép nhóm trên.
	\loigiai{
		Bảng tần số ghép nhóm (theo giá trị đại diện)
		\begin{center}
			\begin{tabular}{|c|c|c|c|c|c|}
				\hline Chiều cao &$[8{,}5; 8{,}8)$ &{$[8{,}8; 9{,}1)$} &{$[9{,}1; 9{,}4)$} &{$[9{,}4; 9{,}7)$} &{$[9{,}7; 10{,}0)$} \\
				\hline Giá trị đại diện &$8{,}65$ &$8{,}95$ &$9{,}25$ &$9{,}55$ &$9{,}85$ \\
				\hline Số cây & $20$ & $35$ & $60$ & $55$ & $30$\\
				\hline
			\end{tabular}
		\end{center}
		Chiều cao trung bình của $200$ cây keo 3 năm tuổi là
		\[\overline{x}=\dfrac{8{,}65\cdot 20+8{,}95\cdot 35+9{,}25\cdot 60+9{,}55\cdot 55+9{,}85\cdot 30}{200}\approx 9{,}31. \]
		Nhóm chứa mốt của mẫu số liệu trên là nhóm $[9{,}1; 9{,}4)$.\\
		Do đó $u_m=9{,}1$; $n_{m-1}=35$; $n_m=60$; $n_{m+1}=55$; $u_{m+1}-u_m=9{,}4-9{,}1=0{,}3$.\\
		Mốt của mẫu số liệu ghép nhóm là
		\[M_0=9{,}1+\dfrac{60-35}{(60-35)+(60-55)} \cdot 0{,}3= 9{,}35. \]
	}
\end{bt}

\begin{bt}%[1D5V2-3]
% nguon: Archive/2024/TOAN_11_XE/TOAN_11_XE/data4/11C3/KNTT-3.2.tex #912

	Một hãng xe ô tô thống kê lại số lần gặp sự cố về động cơ của $ 100 $ chiếc xe cùng loại sau $ 2 $ năm sử dụng đầu tiên ở bảng sau:
	\begin{center}
		\begin{tabular}{|c|c|c|c|c|c|}
			\hline 
			\textbf{Số lần gặp sự cố}	& $ \left[ 1 ; 2\right]  $ & $ \left[3 ; 4\right]  $ & $ \left[ 5 ; 6\right]  $ & $ \left[ 7 ; 8\right]  $ & $ \left[ 9 ; 10\right]  $ \\ 
			\hline 
			\textbf{Số xe}	& $ 17 $  & $ 33 $ & $ 25 $ & $ 20 $ & $ 5 $ \\ 
			\hline 
		\end{tabular} 
	\end{center}
	\begin{enumerate}[a)]
		\item Hãy ước lượng các tứ phân vị của mẫu số liệu ghép nhóm trên.
		\item Một người cho rằng có trên $ 25\% $ xe của hãng gặp không ít hơn $ 4 $ sự cố về đồng cơ trong $ 2 $ năm sử dụng đầu tiên. Nhận định trên có hợp lí không?
	\end{enumerate}
	\loigiai{
		\begin{enumerate}[a)]
			\item Do số lần gặp sự cố là số nguyên nên ta hiệu chỉnh lại như sau:
			\begin{center}
				\begin{tabular}{|c|c|c|c|c|c|}
					\hline 
					\textbf{Số lần gặp sự cố}	& $ \left[ 0{,}5 ; 2{,}5\right)  $ & $ \left[ 2{,}5 ; 4{,}5\right)  $ & $ \left[4{,}5 ; 6{,}5 \right)  $ & $ \left[ 6{,}5 ; 8{,}5\right) $ &$ \left[8{,}5 ; 10{,}5 \right)  $  \\ 
					\hline 
					\textbf{Số xe}	& $ 17 $ & $ 33 $ &$ 25 $  & $ 20 $ & $ 5 $ \\ 
					\hline 
				\end{tabular} 
			\end{center}
			Gọi $ x_1; x_2; \ldots; x_{100} $ là mẫu số liệu được xếp theo thứ tự không giảm.\\
			Ta có $ x_1, \ldots, x_{17}\in \left[ 0{,}5 ; 2{,}5\right)  $;  $ x_{18}, \ldots, x_{50}\in \left[ 2{,}5 ; 4{,}5\right)  $; $ x_{51}, \ldots, x_{75}\in \left[ 4{,}5 ; 6{,}5\right) $; $ x_{76}, \ldots, x_{95}\in \left[ 6{,}5 ; 8{,}5\right)  $; $ x_{96}, \ldots, x_{100}\in \left[ 8{,}5 ; 10{,}5\right)  $.\\
			Tứ phân vị thứ hai của dãy số liệu $ x_1; x_2; \ldots; x_{100} $ là $ \dfrac{1}{2}\left( x_{50}+x_{51}\right) $. Do $ x_{50}\in \left[ 2{,}5 ; 4{,}5\right) $ và $ x_{51}\in \left[4{,}5 ; 6{,}5 \right) $ nên tứ phân vị thứ hai của mẫu số liệu ghép nhóm là $ Q_2=4{,}5$.\\
			Tứ phân vị thứ nhất của dãy số liệu $ x_1; x_2; \ldots; x_{100} $ là $ \dfrac{1}{2}\left( x_{25}+x_{26}\right) $. Do $ x_{25} $ và $ x_{26} $ thuộc nhóm $ \left[2{,}5 ; 4{,}5 \right)$ nên tứ phân vị thứ nhất của mẫu số liệu ghép nhóm là $$ Q_1=2{,}5+\dfrac{\dfrac{1\cdot100}{4}-17}{33}\cdot (4{,}5-2{,}5)=\dfrac{197}{66}\approx 2{,}98.$$
			Tứ phân vị thứ ba của dãy số liệu $ x_1; x_2; \ldots; x_{100} $ là $ \dfrac{1}{2}\left(x_{75}+x_{76} \right)$. Do $ x_{75}\in \left[4{,}5 ; 6{,}5 \right) $ và $ x_{76}\in \left[ 6{,}5 ; 8{,}5\right)$ nên tứ phân vị thứ ba của mẫu số liệu ghép nhóm là $ Q_3=6{,}5$.
			\item Do tứ phân vị thứ nhất $ Q_1\approx 2{,}98 $ nên nhận định trên là hợp lý.
		\end{enumerate}
	}
\end{bt}

\begin{bt}%[1D5V2-3]
% nguon: Archive/2024/TOAN_11_XE/TOAN_11_XE/data4/11C3/KNTT-3.2.tex #913

	Một người thống kê lại thời gian thực hiện các cuộc gọi điện thoại của người đó trong một tuần ở bảng sau:
	\begin{center}
		\begin{tabular}{|c|c|c|c|c|c|c|}
			\hline 
			\begin{tabular}{c}
				\textbf{Thời gian}\\
				\textbf{ (đơn vị: giây)}
			\end{tabular} 
			& $ \left[ 0;60\right)  $ & $ \left[ 60 ; 120\right)  $ & $ \left[ 120 ; 180\right)  $ & $ \left[ 180; 240\right)  $ & $ \left[240 ; 300 \right)  $ & $ \left[300 ; 360 \right)  $ \\ 
			\hline 
			\textbf{Số cuộc gọi}	& $ 8 $ & $ 10 $ & $ 7 $ & $ 5 $ &$ 2 $  & $ 1 $ \\ 
			\hline 
		\end{tabular} 
	\end{center}
	Hãy ước lượng các tứ phân vị của mẫu số liệu ghép nhóm trên.
	\loigiai{
		Gọi $ x_1; x_2; \ldots; x_{33} $ là mẫu số liệu được xếp theo thứ tự không giảm.\\
		Ta có $ x_1, \ldots, x_{8}\in \left[ 0 ; 60\right)  $; $ x_9, \ldots, x_{18} \in \left[ 60 ; 120\right) $, $ x_{19}, \ldots, x_{25} \in \left[ 120 ; 180\right) $;  $ x_{26}, \ldots, x_{30} \in \left[ 180 ; 240\right) $; $ x_{31}, x_{32} \in \left[ 240 ; 300\right) $; $ x_{33}\in \left[ 300 ; 360\right) $.\\
		Tứ phân vị thứ nhất của dãy số liệu $ x_1; x_2; \ldots; x_{33} $ là  $ x_9\in \left[60 ; 120 \right) $ nên tứ phân vị thứ nhất của mẫu số liệu ghép nhóm là $$ Q_1=60+\dfrac{\dfrac{33}{4}-8}{10}\cdot (120-60)=\dfrac{123}{2} \approx 61{,}5.$$
		Tứ phân vị thứ hai của dãy số liệu $ x_1; x_2; \ldots; x_{33} $ là  $ x_{17}\in \left[60 ; 120 \right) $ nên tứ phân vị thứ hai của mẫu số liệu ghép nhóm là $$ Q_2=60+\dfrac{\dfrac{33}{2}-8}{10}\cdot (120-60)=111.$$
		Tứ phân vị thứ ba của dãy số liệu $ x_1; x_2; \ldots; x_{33} $ là  $ x_{25}\in \left[120 ; 180 \right) $ nên tứ phân vị thứ ba của mẫu số liệu ghép nhóm là $$ Q_3=120+\dfrac{\dfrac{33\cdot3}{4}-(8+10)}{7}\cdot (180-120)=\dfrac{1245}{7}\approx 177{,}857.$$
	}
\end{bt}

\begin{bt}%[1D5V2-3]
% nguon: Archive/2024/TOAN_11_XE/TOAN_11_XE/data4/11C3/KNTT-3.2.tex #914

	Một phòng khám thống kê số bệnh nhân đến khám bệnh mỗi ngày trong tháng $ 4 $ năm $ 2022 $ ở bảng sau:
	\begin{center}
		\begin{tabular}{|c|c|c|c|c|c|}
			\hline 
			\textbf{Số bệnh nhân}	& $ \left[1 ; 10 \right]  $ & $ \left[ 11 ; 20\right]  $ & $ \left[ 21 ; 30\right]  $ & $ \left[ 31 ; 40\right]  $ & $ \left[41 ; 50 \right]  $ \\ 
			\hline 
			\textbf{Số ngày}	& $ 7 $ & $ 8 $ & $ 7 $ & $ 6 $ & $ 2 $ \\ 
			\hline 
		\end{tabular} 
	\end{center}
	\begin{enumerate}[a)]
		\item Hãy ước lượng các tứ phân vị của mẫu số liệu ghép nhóm trên.
		\item Quản lý phòng khám cho rằng có khoảng $ 25\% $ số ngày khám có nhiều hơn $ 35 $ bệnh nhân đến khám. Nhận định trên có hợp lý không?
	\end{enumerate}
	\loigiai{
		\begin{enumerate}[a)]
			\item Do số bệnh nhân đến khám là số nguyên nên ta hiệu chỉnh lại như sau:
			\begin{center}
				\begin{tabular}{|c|c|c|c|c|c|}
					\hline 
					\textbf{Số bệnh nhân}	& $ \left[0{,}5 ; 10{,}5 \right)  $ & $ \left[ 10{,}5 ; 20{,}5\right)  $ & $ \left[ 20{,}5 ; 30{,}5\right)  $ & $ \left[ 30{,}5 ; 40{,}5\right)  $ & $ \left[40{,}5 ; 50{,}5 \right)  $ \\ 
					\hline 
					\textbf{Số ngày}	& $ 7 $ & $ 8 $ & $ 7 $ & $ 6 $ & $ 2 $ \\ 
					\hline 
				\end{tabular} 
			\end{center}
			Gọi $ x_1; x_2; \ldots; x_{30} $ là mẫu số liệu được xếp theo thứ tự không giảm.\\
			Ta có $ x_1, \ldots, x_{7}\in \left[ 0{,}5 ; 10{,}5\right)  $; $ x_8, \ldots, x_{15} \in \left[ 10{,}5 ; 20{,}5\right) $, $ x_{16}, \ldots, x_{22} \in \left[ 20{,}5 ; 30{,}5\right) $;  $ x_{23}, \ldots, x_{28} \in \left[ 30{,}5 ; 40{,}5\right) $; $ x_{29}, x_{30} \in \left[ 40{,}5 ; 50{,}5\right) $.\\
			Tứ phân vị thứ nhất của dãy số liệu $ x_1; x_2; \ldots; x_{30} $ là  $ x_8 \in \left[10{,}5 ; 20{,}5 \right) $ nên tứ phân vị thứ nhất của mẫu số liệu ghép nhóm là $$ Q_1=10{,}5+\dfrac{\dfrac{30}{4}-7}{8}\cdot (20{,}5-10{,}5)=\dfrac{89}{8}\approx 11{,}125.$$
			Tứ phân vị thứ hai của dãy số liệu $ x_1; x_2; \ldots; x_{30} $ là $ \dfrac{1}{2}\left( x_{15}+x_{16}\right) $. Do $ x_{15} \in \left[ 10{,}5 ; 20{,}5 \right) $ và $ x_{16}\in \left[ 20{,}5 ; 30{,}5\right) $ nên tứ phân vị thứ hai của mẫu số liệu ghép nhóm là $ Q_2=20{,}5 $.\\
			Tứ phân vị thứ hai của dãy số liệu $ x_1; x_2; \ldots; x_{30} $ là $ x_{23}\in  \left[ 30{,}5 ; 40{,}5\right)$ nên tứ phân vị thứ ba của mẫu số liệu ghép nhóm là $$Q_3= 30{,}5+\dfrac{\dfrac{3\cdot 30}{4}-(7+8+7)}{6}\cdot(40{,}5-30{,}5)=\dfrac{94}{3}\approx 31{,}333.$$
			\item  Do $ Q_1\approx 11{,}125 $ nên nhận định $ 25\% $ số ngày khám nhiều hơn $ 35 $ bệnh nhân là không hợp lý.
		\end{enumerate}
	}
\end{bt}

\begin{bt}%[1D5V2-3]
% nguon: Archive/2024/TOAN_11_XE/TOAN_11_XE/data4/11C3/KNTT-3.2.tex #915

	Lương tháng của một số nhân viên văn phòng được ghi lại như sau (đơn vị: triệu đồng):
	\begin{center}
		\begin{tabular}{|c|c|c|c|c|c|c|c|c|c|c|c|}
			\hline 
			$ 12{,}5 $	& $ 9{,}6 $ & $ 11{,}7 $ & $ 12{,}7 $ & $ 10{,}0 $ & $ 10{,}0 $ & $ 12{,}2 $ & $ 9{,}8 $ & $ 10{,}9 $ & $ 6{,}7 $ & $ 13{,}6 $ & $ 9{,}2 $ \\ 
			\hline 
			$ 13{,}1 $	& $ 6{,}5 $ & $ 10{,}7 $ & $ 8{,}9 $ & $ 11{,}2 $ & $ 13{,}2 $ & $ 8{,}3 $ & $ 11{,}1 $ & $ 11{,}9 $ & $ 8{,}4 $ & $ 6{,}7 $ & $ 13{,}8 $ \\ 
			\hline 
		\end{tabular} 
	\end{center}
	\begin{enumerate}[a)]
		\item Tìm tứ phân vị của dãy số liệu trên.
		\item Tổng hợp lại dãy số liệu trên dựa vào bảng tần số ghép nhóm theo mẫu sau:
		\begin{center}
			\begin{tabular}{|c|c|c|c|c|}
				\hline 
				\begin{tabular}{c}
					\textbf{Lương tháng}
					\\ 
					\textbf{	(triệu đồng)}
				\end{tabular} 
				& $ \left[ 6 ; 8\right)  $ & $ \left[8 ; 10 \right)  $ & $ \left[10 ; 12 \right)  $ & $ \left[ 12 ; 14\right)  $ \\ 
				\hline 
				\textbf{Số nhân viên}& ? & ? & ? & ? \\ 
				\hline 
			\end{tabular} 
		\end{center}
		\item Hãy ước lượng tứ phân vị của số liệu ở bảng tần số ghép nhóm trên.
	\end{enumerate}
	\loigiai{
		\begin{enumerate}[a)]
			\item Gọi $ x_1 $, $ x_2 $, \ldots, $ x_{24} $ là mẫu số liệu được xếp theo thứ tự không giảm. Ta có bảng sau
			\begin{center}
				\begin{tabular}{|c|c|c|c|c|c|c|c|c|c|c|c|c|c|c|c|c|c|c|c|c|c|c|c|}
					\hline 
					$ x_1 $ &$ x_2$  & $ x_3 $ & $ x_4 $ & $ x_5$ &$ x_6$  & $ x_7 $ & $ x_8$ &$ x_9$  & $ x_{10} $ &$ x_{11} $  & $ x_{12} $ \\
					\hline 
					$ 6{,}5 $ & $ 6{,}7 $ & $ 6{,}7 $ & $ 8{,}3 $ & $ 8{,}4 $ & $ 8{,}9 $ & $ 9{,}2 $ & $ 9{,}6 $ & $ 9{,}8 $ & $ 10{,}0 $ &  $ 10{,}0 $ & $ 10{,}7 $\\ 
					\hline 
					$ x_{13} $& $ x_{14} $ & $ x_{15} $ & $ x_{16} $ & $ x_{17} $ & $ x_{18} $ & $ x_{19}$ & $ x_{20} $ & $ x_{21} $ & $ x_{22} $ & $ x_{23} $ & $ x_{24} $ \\ 
					\hline
					$ 10{,}9 $& $ 11{,}1 $ & $ 11{,}2 $ & $ 11{,}7 $ & $ 11{,}9 $ & $ 12{,}2 $ & $ 12{,}5 $ & $ 12{,}7 $ & $ 13{,}1 $ & $ 13{,}2 $ & $ 13{,}6 $ & $ 13{,}8 $  \\ 
					\hline	 
				\end{tabular} 
			\end{center}
			Tứ phân vị thứ nhất của mẫu số liệu là $ \dfrac{1}{2}\cdot\left( x_6+x_7\right)= \dfrac{1}{2}\cdot \left( 8{,}9+9{,}2\right)=\dfrac{181}{10}=18{,}1$.\\
			Tứ phân vị thứ hai của mẫu số liệu là $ \dfrac{1}{2}\cdot\left( x_{12}+x_{13}\right)= \dfrac{1}{2}\cdot \left( 10{,}7+10{,}9\right)=\dfrac{54}{5}=10{,}8$.\\
			Tứ phân vị thứ ba của mẫu số liệu là $\dfrac{1}{2}\cdot\left(  x_{18}+x_{19}\right) = \dfrac{1}{2}\cdot \left( 12{,}2+12{,}5\right) =\dfrac{247}{20}=12{,}35$.
			\item Tổng hợp lại dãy số liệu trên vào bảng tần số ghép nhóm sau:
			\begin{center}
				\begin{tabular}{|c|c|c|c|c|}
					\hline 
					\begin{tabular}{c}
						\textbf{Lương tháng}
						\\ 
						\textbf{	(triệu đồng)}
					\end{tabular} & $ \left[ 6 ; 8\right)  $ & $ \left[8 ; 10 \right)  $ & $ \left[10 ; 12 \right)  $ & $ \left[ 12 ; 14\right)  $ \\ 
					\hline 
					\textbf{Số nhân viên}& $ 3 $ & $ 6 $ & $ 8 $ & $ 7 $ \\ 
					\hline 
				\end{tabular} 
			\end{center}
			\item Gọi $ x_1 $, $ x_2 $, \ldots, $ x_{24} $ là mẫu số liệu được xếp theo thứ tự không giảm. \\
			Ta có $ x_1, \ldots,  x_3 \in \left[ 6 ; 8\right)  $; $ x_4, \ldots, x_9 \in \left[8 ; 10 \right)$; $ x_{10}, \ldots, x_{17}\in \left[10 ; 12 \right) $; $ x_{18}, \ldots, x_{24}\in  \left[ 12 ; 14\right) $.\\
			Tứ phân vị thứ nhất của mẫu số liệu là $ \dfrac{1}{2}\cdot\left( x_6+x_7\right) \in \left[8 ; 10 \right)$ nên tứ phân vị thứ nhất của mẫu số liệu ghép nhóm là $$Q_1= 8+\dfrac{\dfrac{24}{4}-3}{6} \cdot \left( 10-8\right)=9.$$
			Tứ phân vị thứ hai của mẫu số liệu là $ \dfrac{1}{2}\cdot\left( x_{12}+x_{13}\right) \in \left[10 ; 12 \right)$ nên tứ phân vị thứ hai của mẫu số liệu ghép nhóm là $$Q_2= 10+\dfrac{\dfrac{24}{2}-(3+6)}{8} \cdot \left( 12-10\right)=\dfrac{43}{4}=10{,}75.$$
			Tứ phân vị thứ ba của mẫu số liệu là $\dfrac{1}{2}\cdot\left(  x_{18}+x_{19}\right)\in \left[ 12 ; 14\right)$ nên tứ phân vị thứ ba của mẫu số liệu ghép nhóm là $$Q_3= 12+\dfrac{\dfrac{24\cdot 3}{4}-(3+6+8)}{7} \cdot \left( 14-12\right)=\dfrac{86}{7}\approx 12{,}286.$$
		\end{enumerate}
	}
\end{bt}

\begin{bt}%[1D5V2-3]
% nguon: Archive/2024/TOAN_11_XE/TOAN_11_XE/data4/11C3/KNTT-3.2.tex #916

	Số điểm một cầu thủ bóng rổ ghi được trong $ 20 $ trận đấu được cho ở bảng sau:
	\begin{center}
		\begin{tabular}{|c|c|c|c|c|c|c|c|c|c|c|c|c|c|c|c|c|c|c|c|}
			\hline 
			$ 25 $	& $ 23 $ & $ 21 $ & $ 13 $ & $ 8 $ & $ 14 $ & $ 15 $ & $ 18 $ & $ 22 $ & $ 11 $ \\ 
			\hline 
			$ 24 $	& $ 12 $ & $ 14 $ & $ 14 $ & $ 18 $ & $ 6 $ & $ 8 $ & $ 25 $ & $ 10 $ & $ 11 $  \\ 
			\hline 
		\end{tabular} 
	\end{center}
	\begin{enumerate}[a)]
		\item Tìm tứ phân vị của dãy số liệu trên.
		\item Tổng hợp lại dãy số liệu trên vào bảng tần số ghép nhóm theo mẫu sau:
		\begin{center}
			\begin{tabular}{|c|c|c|c|c|}
				\hline 
				\textbf{Điểm số}	& $ \left[ 6 ; 10\right]  $ & $ \left[ 11 ; 15\right]  $ & $ \left[16 ; 20 \right]  $ & $ \left[ 21 ; 25\right]  $ \\ 
				\hline 
				\textbf{Số trận}	& ? & ? & ? & ? \\ 
				\hline 
			\end{tabular} 
		\end{center}
		\item Hãy ước lượng tứ phân vị của số liệu từ bảng tần số ghép nhóm trên.
	\end{enumerate}
	\loigiai{
		\begin{enumerate}[a)]
			\item Gọi $ x_1 $, $ x_2 $, \ldots, $ x_{20} $ là mẫu số liệu được xếp theo thứ tự không giảm. Ta có bảng sau
			\begin{center}
				\begin{tabular}{|c|c|c|c|c|c|c|c|c|c|}
					\hline 
					$ x_1 $	& $ x_2 $ & $ x_3 $ & $ x_4 $ & $ x_5 $ & $ x_6 $ & $ x_7 $ & $ x_8 $ & $ x_9 $ & $ x_{10} $ \\ 
					\hline 
					$ 6 $	& $ 8 $ & $ 8 $ & $ 10 $ & $ 11 $ & $ 11 $ & $ 12 $ & $ 13 $ & $ 14 $ & $ 14 $ \\ 
					\hline 
					$ x_{11} $	& $ x_{12} $ & $ x_{13} $ & $ x_{14} $ & $ x_{15} $ & $ x_{16} $ & $ x_{17} $ & $ x_{18} $ & $ x_{19} $ & $ x_{20} $ \\ 
					\hline 
					$ 14 $	& $ 15 $ & $ 18 $ & $ 18 $ & $ 21 $ & $ 22 $ & $ 23 $ & $ 24 $ & $ 25 $ & $ 25 $ \\
					\hline 
				\end{tabular} 
			\end{center}
			Tứ phân vị thứ nhất của mẫu số liệu là $ \dfrac{1}{2}\cdot\left( x_{5}+x_6\right) =\dfrac{1}{2}\cdot \left( 11+11\right)=11$.\\
			Tứ phân vị thứ hai của mẫu số liệu là $ \dfrac{1}{2}\cdot \left( x_{10}+x_{11}\right)= \dfrac{1}{2}\cdot \left( 14+14\right)=14$.\\
			Tứ phân vị thứ ba của mẫu số liệu là $ \dfrac{1}{2}\cdot \left( x_{15}+x_{16}\right)= \dfrac{1}{2}\cdot \left( 21+22\right)=21{,}5$.
			\item Tổng hợp lại dãy số liệu trên vào bảng tần số ghép nhóm sau:
			\begin{center}
				\begin{tabular}{|c|c|c|c|c|}
					\hline 
					\textbf{Điểm số}	& $ \left[ 6 ; 10\right]  $ & $ \left[ 11 ; 15\right]  $ & $ \left[16 ; 20 \right]  $ & $ \left[ 21 ; 25\right]  $ \\ 
					\hline 
					\textbf{Số trận}	& $ 4 $ & $ 8 $ & $ 2 $ & $ 6 $ \\ 
					\hline 
				\end{tabular} 
			\end{center}
			\item Hãy ước lượng tứ phân vị của số liệu từ bảng tần số ghép nhóm trên.\\
			Vì số điểm của cầu thủ bóng rổ là số nguyên nên ta hiệu chỉnh bảng số liệu  tần số ghép nhóm như sau:
			\begin{center}
				\begin{tabular}{|c|c|c|c|c|}
					\hline 
					\textbf{Điểm số}	& $ \left[ 5{,}5 ; 10{,}5\right)  $ & $ \left[ 10{,}5 ; 15{,}5\right)  $ & $ \left[15{,}5 ; 20{,}5 \right)  $ & $ \left[ 20{,}5 ; 25{,}5\right)  $ \\ 
					\hline 
					\textbf{Số trận}	& $ 4 $ & $ 8 $ & $ 2 $ & $ 6 $ \\ 
					\hline 
				\end{tabular} 
			\end{center}
			Gọi $ x_1 $, $ x_2 $, $ \ldots $, $ x_{20} $ là mẫu số liệu được xếp theo thứ tự không giảm.\\
			Ta có $ x_1, \ldots, x_4 \in  \left[ 5{,}5 ; 10{,}5\right) $; $ x_5 \ldots, x_{12} \in   \left[ 10{,}5 ; 15{,}5\right)$; $ x_{13}, x_{14}\in  \left[15{,}5 ; 20{,}5 \right) $; $ x_{15}, \ldots, x_{20} \in  \left[ 20{,}5 ; 25{,}5\right) $.\\
			Tứ phân vị thứ nhất của mẫu số liệu là $ \dfrac{1}{2}\cdot\left( x_{5}+x_6\right) \in \left[ 10{,}5 ; 15{,}5\right)$ nên tứ phân vị thứ nhất của mẫu số liệu ghép nhóm là $$Q_1=10{,}5 +\dfrac{\dfrac{20}{4}-4}{8}\cdot \left(15{,}5-10{,}5 \right)=\dfrac{89}{8}=11{,}125.$$
			Tứ phân vị thứ hai của mẫu số liệu là $ \dfrac{1}{2}\cdot\left( x_{10}+x_{11}\right) \in \left[ 10{,}5 ; 15{,}5\right)$ nên tứ phân vị thứ hai của mẫu số liệu ghép nhóm là $$Q_2=10{,}5 +\dfrac{\dfrac{20}{2}-4}{8}\cdot \left(15{,}5-10{,}5 \right)=\dfrac{57}{4}=14{,}25.$$
			Tứ phân vị thứ ba của mẫu số liệu là $ \dfrac{1}{2}\cdot\left( x_{15}+x_{16}\right) \in \left[ 20{,}5 ; 25{,}5\right)$ nên tứ phân vị thứ ba của mẫu số liệu ghép nhóm là $$Q_3=20{,}5 +\dfrac{\dfrac{20\cdot 3}{4}-(4+8+2)}{6}\cdot \left(25{,}5-20{,}5 \right)=\dfrac{64}{3}\approx 21{,}333.$$
		\end{enumerate}
	}
\end{bt}

\begin{bt}%[1D5V2-3]
% nguon: Archive/2024/TOAN_11_XE/TOAN_11_XE/data4/11C3/KNTT-3.2.tex #917

	Kiểm tra điện lượng của một số viên pin tiểu do một hãng sản xuất thu được kết quả như sau:
	\begin{center}
		\begin{tabular}{|c|c|c|c|c|c|}
			\hline 
			\begin{tabular}{c}
				\textbf{Điện lượng} \\	\textbf{(nghìn mAh)}
			\end{tabular} 
			& $ \left[ 0{,}9 ; 0{,}95\right)  $ & $ \left[ 0{,}95 ; 1{,}0\right)  $ & $ \left[ 1{,}0 ; 1{,}05\right)  $ &$ \left[ 1{,}05 ; 1{,}1\right)  $  &  $ \left[ 1{,}1 ; 1{,}15\right)  $\\ 
			\hline 
			\textbf{Số viên pin}& $ 10 $ & $ 20 $ & $ 35 $ & $ 15 $ & $ 5 $ \\ 
			\hline 
		\end{tabular} 
	\end{center}
	Hãy ước lượng số trung bình, mốt và tứ phân vị của mẫu số liệu ghép nhóm trên.
	\loigiai{
		\begin{itemize}
			\item Tìm số trung bình của mẫu số liệu ghép nhóm.\\
			Ta có bảng thống kê điện lượng của pin theo giá trị đại diện là:
			\begin{center}
				\begin{tabular}{|c|c|c|c|c|c|}
					\hline 
					\begin{tabular}{c}
						\textbf{Điện lượng} \\	\textbf{(nghìn mAh)}
					\end{tabular} 
					& $ \left[ 0{,}9 ; 0{,}95\right)  $ & $ \left[ 0{,}95 ; 1{,}0\right)  $ & $ \left[ 1{,}0 ; 1{,}05\right)  $ &$ \left[ 1{,}05 ; 1{,}1\right)  $  &  $ \left[ 1{,}1 ; 1{,}15\right)  $\\ 
					\hline 
					\textbf{Giá trị đại diện}& $ 0{,}925 $ & $ 0{,}975 $ & $ 1{,}025 $ & $ 1{,}075 $ & $ 1{,}125 $ \\ 
					\hline
					\textbf{Số viên pin}& $ 10 $ & $ 20 $ & $ 35 $ & $ 15 $ & $ 5 $ \\ 
					\hline 
				\end{tabular} 
			\end{center}
			Số trung bình của mẫu số liệu ghép nhóm theo dõi điện lượng của một số viên pin xấp xỉ bằng: $$ \dfrac{0{,}925\cdot 10 + 0{,}975\cdot 20 +1{,}025 \cdot 35 +1{,}075 \cdot 15+1{,}125 \cdot 5}{10+20+35+15+5}\approx 1{,}016. $$
			\item Tìm mốt của mẫu số liệu ghép nhóm.\\
			Nhóm chứa mốt của mẫu số liệu là $ \left[ 1{,}0 ; 1{,}05\right) $.\\
			Do đó $ u_m=1$; $ n_{m-1}=20 $; $ n_m=35 $; $ n_{m+1}=15 $; $ u_{m+1}-u_m=1{,}05-1{,}0=0{,}05 $.\\
			Mốt của mẫu số liệu ghép nhóm là $$M_0=1+\dfrac{35-20}{(35-20)+(35-15)}\cdot 0{,}05 =\dfrac{143}{140}\approx 1{,}021.$$  
			\item Tìm các tứ phân vị của mẫu số liệu ghép nhóm.\\
			Gọi $ x_1 $, $ x_2 $, $ \ldots $, $ x_{85} $ là mẫu số liệu được xếp theo thứ tự không giảm.\\
			Ta có $ x_1, \ldots, x_{10}\in  \left[ 0{,}9 ; 0{,}95\right)$; $ x_{11}, \ldots, x_{30}\in  \left[ 0{,}95 ; 1{,}0\right) $; $ x_{31}, \ldots, x_{65}\in   \left[ 1{,}0 ; 1{,}05\right) $; $ x_{66}, \ldots, x_{80}\in \left[ 1{,}05 ; 1{,}1\right) $; $ x_{81}, \ldots, x_{85}\in \left[ 1{,}1 ; 1{,}15\right)$.\\
			Tứ phân vị thứ nhất của mẫu số liệu là $ x_{22} \in  \left[ 0{,}95 ; 1{,}0\right) $ nên tứ phân vị thứ nhất của mẫu số liệu ghép nhóm là 
			$$Q_1=0{,}95+\dfrac{\dfrac{85}{4}-10}{20}\cdot \left( 1{,}0-0{,}95\right)=  \dfrac{313}{320}\approx 0{,}978.$$
			Tứ phân vị thứ hai của mẫu số liệu là $ x_{43} \in  \left[ 1{,}0 ; 1{,}05\right) $ nên tứ phân vị thứ hai của mẫu số liệu ghép nhóm là 
			$$Q_2=1{,}0+\dfrac{\dfrac{85}{2}-(10+20)}{35}\cdot \left( 1{,}05-1{,}0\right)=\dfrac{57}{56} \approx 1{,}018.$$
			Tứ phân vị thứ ba của mẫu số liệu là $ x_{64} \in  \left[ 1{,}0 ; 1{,}05\right) $ nên tứ phân vị thứ ba của mẫu số liệu ghép nhóm là 
			$$Q_3=1{,}0+\dfrac{\dfrac{85\cdot 3}{4}-(10+20)}{35}\cdot \left( 1{,}05-1{,}0\right)=\dfrac{587}{560} \approx 1{,}048.$$
		\end{itemize}
	}
\end{bt}

\begin{bt}%[1D5C2-3]
% nguon: Archive/2024/TOAN_11_XE/TOAN_11_XE/data4/11C3/KNTT-3.2.tex #918

	Cân nặng của một số lợn con mới sinh thuộc hai giống $ A $ và $ B $ được cho ở biểu đồ dưới đây (đơn vị: kg).
	\begin{center}
		\begin{tikzpicture}[>=stealth,line join=round,line cap=round,font=\footnotesize,scale=0.85,line width=1pt]
			\draw[->] (0,0)--(0,5)node[left]{(\text{Số con})};
			\foreach \y in {1,2,3,4}
			\draw[shift={(0,\y)}] (0,0)--(-2pt,0) node[left]{\scriptsize ${\y}0$};
			%	\path (4.5,6) node {\normalsize{\textbf{Cân nặng của một số lợn con mới sinh}}};
			\path (4.5,5.5) node {
				$\begin{array}{c}
					\normalsize{\textbf{Cân nặng của một số}}\\
					\normalsize{\textbf{lợn con mới sinh}}
				\end{array}$
			};
			%% nhãn
			\path (2.5,-1.5) node[rectangle,fill=cyan,draw=none]{};
			\path (3.6,-1.5) node {\text{Giống $ A $}};
			\path (5,-1.5) node[rectangle,fill=orange,draw=none]{};
			\path (6.1,-1.5) node {\text{Giống $ B $}};
			% đường gióng
			\foreach \y in {1,2,3,4}{
				\draw[line width=0.2pt] (0,\y)--(8.4,\y);
			}
			%% cột
			\draw[fill=cyan,draw=none] (0,0)--(0,0.8)--(1,0.8)node[midway,above]{$ 8 $}--(1,0)--cycle;
			\draw[fill=orange,draw=none] (1,0)--(1,1.3)--(2,1.3)node[midway,above]{$ 13 $}--(2,0)--cycle;
			\draw[fill=cyan,draw=none] (2,0)--(2,2.8)--(3,2.8)node[midway,above]{$ 28 $}--(3,0)--cycle;
			\draw[fill=orange,draw=none] (3,0)--(3,1.4)--(4,1.4)node[midway,above]{$ 14 $}--(4,0)--cycle;
			\draw[fill=cyan,draw=none] (4,0)--(4,3.2)--(5,3.2)node[midway,above]{$ 32 $}--(5,0)--cycle;
			\draw[fill=orange,draw=none] (5,0)--(5,2.4)--(6,2.4)node[midway,above]{$ 24 $}--(6,0)--cycle;
			\draw[fill=cyan,draw=none] (6,0)--(6,1.7)--(7,1.7)node[midway,above]{$ 17 $}--(7,0)--cycle;
			\draw[fill=orange,draw=none] (7,0)--(7,1.4)--(8,1.4)node[midway,above]{$ 14 $}--(8,0)--cycle;
			%% miền
			\node [below] at (1,0){$ \left[1{,}0 ; 1{,}1 \right)$};
			\node [below] at (3,0){$ \left[1{,}1 ; 1{,}2 \right)$};
			\node [below] at (5,0){$ \left[1{,}2 ; 1{,}3 \right)$};
			\node [below] at (7,0){$ \left[1{,}3 ; 1{,}4 \right)$};
			\draw[->] (0,0)node [below left=-2pt]{$ O $}--(9,0)node[below]{(\text{kg})};
		\end{tikzpicture}
	\end{center}
	\begin{enumerate}[a)]
		\item Hãy so sánh cân nặng của lợn con mới sinh giống $ A $ và giống $ B $ theo số trung bình và trung vị.
		\item Hãy ước lượng tứ phân vị thứ nhất và thứ ba của cân nặng lợn con mới sinh giống $ A $ và của cân nặng lợn con mới sinh giống $ B $.
	\end{enumerate}
	\loigiai{
		\begin{enumerate}[a)]
			\item Bảng tần số ghép nhóm thống kê cân nặng của lợn con mới sinh giống $ A $ và giống $ B $ như sau:
			\begin{center}
				\begin{tabular}{|c|c|c|c|c|}
					\hline 
					\textbf{Cân nặng (kg)}	& $ \left[1{,}0 ; 1{,}1 \right)$  &$ \left[1{,}1 ; 1{,}2 \right)$  &$ \left[1{,}2 ; 1{,}3 \right)$  & $ \left[1{,}3 ; 1{,}4 \right)$ \\ 
					\hline 
					\textbf{Giá trị đại diện (kg)}	& $1{,}05 $ & $ 1{,}15 $ & $ 1{,}25 $ & $ 1{,}35 $ \\ 
					\hline 
					\begin{tabular}{c}
						\textbf{Giống A}
						\\ 
						\textbf{(đơn vị: con)}
					\end{tabular} 	& $ 8 $ & $ 28 $ & $ 32 $ & $ 17 $ \\ 
					\hline 
					\begin{tabular}{c}
						\textbf{Giống B}
						\\ 
						\textbf{(đơn vị: con)}
					\end{tabular} 	& $ 13 $ & $ 14 $ & $ 24 $ & $ 14 $ \\ 
					\hline 
				\end{tabular} 
			\end{center}
			Cân nặng trung bình của lợn con mới sinh giống $ A $ là $$ \dfrac{1{,}05\cdot 8 + 1{,}15 \cdot 28 + 1{,}25 \cdot 32 + 1{,}35 \cdot 17}{8+28+32+17}=\dfrac{2071}{1700}\approx 1{,}218.$$
			Cân nặng trung bình của lợn con mới sinh giống $ B $ là $$ \dfrac{1{,}05\cdot 13 + 1{,}15 \cdot 14 + 1{,}25 \cdot 24 + 1{,}35 \cdot 14}{13+14+24+14}=\dfrac{121}{100} \approx 1{,}21.$$
			Suy ra cân nặng trung bình của lợn con mới sinh giống $ A $  lớn hơn cân nặng trung bình của lợn con mới sinh giống $ B $. \\
			Trung vị của mẫu số liệu ghép nhóm cân nặng của lợn con  giống $ A $ là $$M_e=1{,}2+\dfrac{\dfrac{85}{2}-(8+28)}{32}\cdot (1{,}3-1{,}2)=\dfrac{781}{640}\approx 1{,}22.$$
			Trung vị của mẫu số liệu ghép nhóm cân nặng của lợn con  giống  $ B $ là $$M_e=1{,}2+\dfrac{\dfrac{65}{2}-(13+14)}{24}\cdot (1{,}3-1{,}2)=\dfrac{587}{480}\approx 1{,}22.$$
			Suy ra trung vị của mẫu số liệu ghép nhóm cân nặng của của lợn con  giống  $ A $ bằng trung vị của mẫu số liệu ghép nhóm cân nặng của của lợn con  giống $ B $.
			\item Hãy ước lượng tứ phân vị thứ nhất và thứ ba của cân nặng lợn con mới sinh giống $ A $ và của cân nặng lợn con mới sinh giống $ B $.
			\begin{itemize}
				\item Đối với mẫu dữ liệu ghép nhóm cân nặng lợn con mới sinh giống $ A $ có:\\
				Tứ phân vị thứ nhất của mẫu số liệu ghép nhóm cân nặng lợn con mới sinh giống $ A $ thuộc $ \left[ 1{,}1 ; 1{,}2\right)  $ nên $$ Q_1= 1{,}1+\dfrac{\dfrac{85}{4}-8}{28}\cdot \left( 1{,}2-1{,}1\right)= \dfrac{257}{224}\approx 1{,}147.$$
				Tứ phân vị thứ ba của mẫu số liệu ghép nhóm cân nặng lợn con mới sinh giống $ A $ thuộc $ \left[ 1{,}2 ; 1{,}3\right)  $ nên $$ Q_3= 1{,}2+\dfrac{\dfrac{85\cdot3 }{4}-(8+28)}{32}\cdot \left( 1{,}3-1{,}2\right)= \dfrac{1647}{1280}\approx 1{,}287.$$
				\item Đối với mẫu dữ liệu ghép nhóm cân nặng lợn con mới sinh giống $ B $ có:\\
				Tứ phân vị thứ nhất của mẫu số liệu ghép nhóm cân nặng lợn con mới sinh giống $ B $ thuộc $ \left[ 1{,}1 ; 1{,}2\right)  $ nên $$ Q_1= 1{,}1+\dfrac{\dfrac{65}{4}-13}{14}\cdot \left( 1{,}2-1{,}1\right)= \dfrac{629}{560}\approx 1{,}123.$$	
				Tứ phân vị thứ ba của mẫu số liệu ghép nhóm cân nặng lợn con mới sinh giống $ B $ thuộc $ \left[ 1{,}2 ; 1{,}3\right)  $ nên $$ Q_3= 1{,}2+\dfrac{\dfrac{65\cdot 3}{4}-(13+14)}{24}\cdot \left( 1{,}3-1{,}2\right)= \dfrac{413}{320}\approx 1{,}291.$$	
			\end{itemize}
		\end{enumerate}		
	}
\end{bt}

\begin{bt}%[1D5H2-2]
% nguon: Archive/2024/TOAN_11_XE/TOAN_11_XE/data4/11C3/KNTT-3.2.tex #920
%[Tex hóa SGK CD, TVN-223]
	\immini
	{
		Sau khi kiểm tra về số học sinh trong $100$ lớp học, người ta chia mẫu số liệu đó thành năm nhóm căn cứ vào số lượng học sinh của mỗi lớp (đơn vị: học sinh) và lập bảng tần số ghép nhóm bao gồm tần số tích luỹ như bảng bên. Tìm trung vị của mẫu số liệu đó.
	}
	{
		\begin{tabular}{|c|c|c|}
			\hline
			\textbf{Nhóm} & \textbf{Tần số} & \textbf{Tần số tích luỹ}\\ 
			\hline
			$\left[36;38\right)$ & $9$ & $9$\\
			$\left[38;40\right)$ & $15$ & $24$\\
			$\left[40;42\right)$ & $25$ & $49$\\
			$\left[42;44\right)$ & $30$ & $79$\\
			$\left[44;46\right)$ & $21$ & $100$\\
			\hline
			& $n = 100$ &\\
			\hline
		\end{tabular}
	}
	\loigiai{
		Số phần tử của mẫu là $n=100$. Ta có $\dfrac{n}{2} = \dfrac{100}{2} = 50$.\\
		Do $cf_3 = 49 < 50 < cf_4 = 79$ nên nhóm $4$ là nhóm đầu tiên có tần số tích luỹ lớn hơn hoặc bằng $50$.\\
		Xét nhóm $4$ là nhóm $\left[42;44\right)$ có $r=42$; $d=2$ và $n_4=30$ và nhóm $3$ là nhóm $\left[40;42\right)$ có $cf_3 = 49$.\\
		Khi đó trung vị của mẫu số liệu là 
		\[
		M_e = 42 + \dfrac{50 - 49}{30} \cdot 2 \approx 42\text{ (học sinh)}.
		\]
	}
\end{bt}

\begin{bt}%[1D5H2-3]
% nguon: Archive/2024/TOAN_11_XE/TOAN_11_XE/data4/11C3/KNTT-3.2.tex #921
%[Tex hóa SGK CD, TVN-223]
	Bảng bên cho biết tần số ghép nhóm số liệu thống kê cân nặng của $40$ học sinh lớp $11A$ trong một trường trung học phổ thông (đơn vị: kilôgam). Xác định tứ phân vị của mẫu số liệu ghép nhóm.
	\begin{center}
		\begin{tabular}{|c|c|c|}
			\hline
			\textbf{Nhóm} & \textbf{Tần số} & \textbf{Tần số tích luỹ}\\ 
			\hline
			$\left[30;40\right)$ & $2$ & $2$\\
			$\left[40;50\right)$ & $10$ & $12$\\
			$\left[50;60\right)$ & $16$ & $28$\\
			$\left[60;70\right)$ & $8$ & $36$\\
			$\left[70;80\right)$ & $2$ & $38$\\
			$\left[80;90\right)$ & $2$ & $40$\\
			\hline
			& $n = 40$ &\\
			\hline
		\end{tabular}
	\end{center}
	\loigiai{
		Số phần tử của mẫu là $n=40$.
		\begin{itemize}
			\item Ta có $\dfrac{n}{4} = 10$, mà $2<10<12$ nên nhóm $2$ là nhóm đầu tiên có tần số tích luỹ lớn hơn hoặc bằng $10$. \\
			Xét nhóm $2$ là nhóm $\left[40;50\right)$ có $s=40$, $h=10$, $n_2=10$ và nhóm $1$ có $cf_1 = 2$.\\
			Khi đó, tứ phân vị thứ nhất là
			\[
			Q_1 = 40 + \dfrac{10-2}{10} \cdot 10 = 48\text{ (kg)}.
			\]
			\item Ta có $\dfrac{n}{2} = 20$, mà $12<20<28$ nên nhóm $3$ là nhóm đầu tiên có tần số tích luỹ lớn hơn hoặc bằng $20$. \\
			Xét nhóm $3$ là nhóm $\left[50;60\right)$ có $r=50$, $d=10$, $n_3=16$ và nhóm $2$ có $cf_2 = 12$.\\
			Khi đó, tứ phân vị thứ hai là
			\[
			Q_2 = 50 + \dfrac{20-12}{16} \cdot 10 = 55\text{ (kg)}.
			\]
			\item Ta có $\dfrac{3n}{4} = 30$, mà $28<30<36$ nên nhóm $4$ là nhóm đầu tiên có tần số tích luỹ lớn hơn hoặc bằng $30$. \\
			Xét nhóm $4$ là nhóm $\left[60;70\right)$ có $t=50$, $l=10$, $n_4=8$ và nhóm $3$ có $cf_3 = 28$.\\
			Khi đó, tứ phân vị thứ ba là
			\[
			Q_3 = 60 + \dfrac{30-28}{8} \cdot 10 = 62{,}5\text{ (kg)}.
			\]
		\end{itemize}
		Vậy tứ phân vị của mẫu số liệu trên là $48$, $55$ và $62{,}5$.
	}
\end{bt}

\begin{bt}%[1D5H1-4]
% nguon: Archive/2024/TOAN_11_XE/TOAN_11_XE/data4/11C3/KNTT-3.2.tex #922
%[Tex hóa SGK CD, TVN-223]
	Kết quả kiểm tra môn Toán của lớp $11D$ như sau
	\[
	\begin{array}{cccccccccccccccccccc}
		5 & 6 & 7 & 5 & 6 & 9 & 10 & 8 & 5 & 5 & 4 & 5 & 4 & 5 & 7 & 4 & 5 & 8 & 9 & 10 \\
		5 & 3 & 5 & 6 & 5 & 7 & 5 & 8 & 4 & 9 & 5 & 6 & 5 & 6 & 8 & 8 & 7 & 9 & 7 & 9
	\end{array}
	\]
	\begin{enumerate}
		\item Lập bảng tần số ghép nhóm của mẫu số liệu trên có bốn nhóm ứng với bốn nửa khoảng $\left[3;5\right)$, $\left[5;7\right)$, $\left[7;9\right)$, $\left[9;11\right)$.
		\item Mốt của bảng số liệu ghép nhóm trên là bao nhiêu (làm tròn kết quả đến hàng phần mười)?
	\end{enumerate}
	\loigiai{
		\immini
		{
			\begin{enumerate}
				\item Bảng bên là bảng tần số ghép nhóm cho kết quả kiểm tra môn Toán của lớp $11D$.
				\item Ta thấy: Nhóm $2$ ứng với nửa khoảng $\left[5;7\right)$ là nhóm có tần số lớn nhất với $u=5$, $g=2$, $n_2 = 18$. Nhóm $1$ có tần số $n_1 = 5$, nhóm $3$ có tần số $n_3=10$.\\
				Khi đó, mốt của mẫu số liệu là 
				\[
				M_o = 5 + \left( \dfrac{18- 5}{2\cdot 18 - 5 - 10} \right) \cdot 2 \approx 6{,}2.
				\]
			\end{enumerate}
		}
		{
			\begin{tabular}{|c|c|}
				\hline
				\textbf{Nhóm} & \textbf{Tần số}\\ 
				\hline
				$\left[3;5\right)$ & $5$\\
				\hline
				$\left[5;7\right)$ & $18$\\
				\hline
				$\left[7;9\right)$ & $10$\\
				\hline
				$\left[9;11\right)$ & $7$\\
				\hline
				& $n = 40$ \\
				\hline
			\end{tabular}
		}
	}
\end{bt}

\begin{bt}%[1D5H2-3]
% nguon: Archive/2024/TOAN_11_XE/TOAN_11_XE/data4/11C3/KNTT-3.2.tex #923
%[Tex hóa SGK CD, TVN-223]
	Mẫu số liệu dưới đây ghi lại tốc độ của $40$ ô tô khi đi qua một trạm đo tốc độ (đơn vị: km/h):
	\[
	\begin{array}{cccccccccc}
		48{,}5 & 43 & 50 & 55 & 45 & 60 & 53 & 55,5 & 44 & 65 \\
		51 & 62,5 & 41 & 44,5 & 57 & 57 & 68 & 49 & 46{,}5 & 53{,}5 \\
		61 & 49{,}5 & 54 & 62 & 59 & 56 & 47 & 50 & 60 & 61 \\
		49{,}5 & 52{,}5 & 57 & 47 & 60 & 55 & 45 & 47,5 & 48 & 61{,}5
	\end{array}
	\]
	\begin{enumerate}
		\item Lập bảng tần số ghép nhóm cho mẫu số liệu trên có sáu nhóm ứng với sáu nửa khoảng:
		\[
		[40 ; 45),[45 ; 50),[50 ; 55),[55 ; 60),[60 ; 65),[65 ; 70).
		\]
		\item Xác định số trung bình cộng, trung vị, tứ phân vị của mẫu số liệu ghép nhóm trên.
		\item Mốt của mẫu số liệu ghép nhóm trên là bao nhiêu?
	\end{enumerate}
	\loigiai{
		\begin{enumerate}
			\item Ta có bảng tần số ghép nhóm của mẫu số liệu trên như sau:
			\begin{center}
				\begin{tabular}{|c|c|c|c|}
					\hline
					\textbf{Nhóm} & \textbf{Giá trị đại diện} & \textbf{Tần số} & \textbf{Tần số tích luỹ}\\ 
					\hline
					$\left[40;45\right)$ & $42{,}5$ & $4$ & $4$\\
					$\left[45;50\right)$ & $47{,}5$ & $11$ & $15$\\
					$\left[50;55\right)$ & $52{,}5$ & $7$ & $22$\\
					$\left[55;60\right)$ & $57{,}5$ & $8$ & $30$\\
					$\left[60;65\right)$ & $62{,}5$ & $8$ & $38$\\
					$\left[65;70\right)$ & $67{,}5$ & $2$ & $40$\\
					\hline
					&  & $n = 40$ &\\
					\hline
				\end{tabular}
			\end{center}
			\item 
			\begin{itemize}
				\item Trung bình cộng của mẫu số liệu trên là
				\[
				\overline{x} = \dfrac{42{,}5 \cdot 4 + 47{,}5 \cdot 11 + 52{,}5 \cdot 7+ 57{,}5 \cdot 8+ 62{,}5 \cdot 8 + 67{,}5 \cdot 2}{40} = 53{,}875\text{ (km/h)}.
				\]
				\item Ta có $\dfrac{n}{2} = 20$, mà $15<20<22$ nên nhóm $3$ là nhóm đầu tiên có tần số tích luỹ lớn hơn hoặc bằng $20$. \\
				Xét nhóm $3$ là nhóm $\left[50;55\right)$ có $r=50$, $d=5$, $n_3=7$ và nhóm $2$ có $cf_2 = 15$.\\
				Khi đó, tứ phân vị thứ hai (cũng là trung vị) là
				\[
				Q_2 = M_e = 50 + \dfrac{20-15}{7} \cdot 5 = 53{,}6\text{ (km/h)}.
				\]
				\item Ta có $\dfrac{n}{4} = 10$, mà $4<10<15$ nên nhóm $2$ là nhóm đầu tiên có tần số tích luỹ lớn hơn hoặc bằng $10$. \\
				Xét nhóm $2$ là nhóm $\left[45;50\right)$ có $s=45$, $h=5$, $n_2=11$ và nhóm $1$ có $cf_1 = 4$.\\
				Khi đó, tứ phân vị thứ nhất là
				\[
				Q_1 = 45 + \dfrac{10-4}{11} \cdot 5 \approx 47{,}7\text{ (km/h)}.
				\]
				\item Ta có $\dfrac{3n}{4} = 30$, mà $30\le 30<38$ nên nhóm $4$ là nhóm đầu tiên có tần số tích luỹ lớn hơn hoặc bằng $30$. \\
				Xét nhóm $4$ là nhóm $\left[55;60\right)$ có $t=55$, $l=5$, $n_4=8$ và nhóm $3$ có $cf_3 = 22$.\\
				Khi đó, tứ phân vị thứ ba là
				\[
				Q_3 = 55 + \dfrac{30-22}{8} \cdot 5 = 60\text{ (km/h)}.
				\]
			\end{itemize}
			Vậy $53{,}875\text{ (km/h)}$ là trung bình mẫu, $53{,}6\text{ (km/h)}$ là trung vị của mẫu và $Q_1 = 47{,}7\text{ (km/h)}$, $Q_2 = 53{,}6\text{ (km/h)}$, $Q_3 = 60\text{ (km/h)}$ là tứ phân vị của mẫu.
			\item Ta thấy: Nhóm $2$ ứng với nửa khoảng $\left[45;50\right)$ là nhóm có tần số lớn nhất với $u=45$, $g=5$, $n_2 = 11$. Nhóm $1$ có tần số $n_1 = 4$, nhóm $3$ có tần số $n_3 = 7$.\\
			Khi đó, mốt của mẫu số liệu là 
			\[
			M_o = 45 + \left( \dfrac{11 - 4}{2\cdot 11 - 4 - 7} \right) \cdot 5 \approx 48{,}2\text{ (km/h)}.
			\]
		\end{enumerate}
	}
\end{bt}

\begin{bt}%[1D5H2-3]
% nguon: Archive/2024/TOAN_11_XE/TOAN_11_XE/data4/11C3/KNTT-3.2.tex #924
%[Tex hóa SGK CD, TVN-223]
	Mẫu số liệu sau ghi lại cân nặng của $30$ bạn học sinh (đơn vị: kilôgam):
	\[
	\begin{array}{cccccccccc}
		17 & 40 & 39 & 40{,}5 & 42 & 51 & 41{,}5 & 39 & 41 & 30\\
		40 & 42 & 40{,}5 & 39{,}5 & 41 & 40{,}5 & 37 & 39{,}5 & 40 & 41\\
		38{,}5 & 39{,}5 & 40 & 41 & 39 & 40{,}5 & 40 & 38{,}5 & 39{,}5 & 41{,}5
	\end{array}
	\]
	\begin{enumerate}
		\item Lập bảng tần số ghép nhóm cho mẫu số liệu trên có tám nhóm ứng với tám nửa khoảng:
		\[
		[15 ; 20),[20 ; 25),[25 ; 30),[30 ; 35),[35 ; 40),[40 ; 45),[45 ; 50),[50 ; 55).
		\]
		\item Xác định số trung bình cộng, trung vị, tứ phân vị của mẫu số liệu ghép nhóm trên.
		\item Mốt của mẫu số liệu ghép nhóm trên là bao nhiêu?
	\end{enumerate}
	\loigiai{
		\begin{enumerate}
			\item Ta có bảng tần số ghép nhóm của mẫu số liệu trên như sau:
			\begin{center}
				\begin{tabular}{|c|c|c|c|}
					\hline
					\textbf{Nhóm} & \textbf{Giá trị đại diện} & \textbf{Tần số} & \textbf{Tần số tích luỹ}\\ 
					\hline
					$\left[15;20\right)$ & $17{,}5$ & $1$ & $1$\\
					$\left[20;25\right)$ & $22{,}5$ & $0$ & $1$\\
					$\left[25;30\right)$ & $27{,}5$ & $0$ & $1$\\
					$\left[30;35\right)$ & $32{,}5$ & $1$ & $2$\\
					$\left[35;40\right)$ & $37{,}5$ & $10$ & $12$\\
					$\left[40;45\right)$ & $42{,}5$ & $17$ & $29$\\
					$\left[45;50\right)$ & $47{,}5$ & $0$ & $29$\\
					$\left[50;55\right)$ & $52{,}5$ & $1$ & $30$\\
					\hline
					&  & $n = 30$ &\\
					\hline
				\end{tabular}
			\end{center}
			\item 
			\begin{itemize}
				\item Trung bình cộng của mẫu số liệu trên là
				\[
				\overline{x} = \dfrac{17{,}5 \cdot 1 + 32{,}5 \cdot 1 + 37{,}5 \cdot 10+ 42{,}5 \cdot 17+ 52{,}5 \cdot 1}{30} = 	40\text{ (kg)}.
				\]
				\item Ta có $\dfrac{n}{2} = 15$, mà $12<15<29$ nên nhóm $6$ là nhóm đầu tiên có tần số tích luỹ lớn hơn hoặc bằng $15$. \\
				Xét nhóm $6$ là nhóm $\left[40;45\right)$ có $r=40$, $d=5$, $n_6=17$ và nhóm $5$ có $cf_5 = 12$.\\
				Khi đó, tứ phân vị thứ hai (cũng là trung vị) là
				\[
				Q_2 = M_e = 40 + \dfrac{15-12}{17} \cdot 5 = 40{,}9\text{ (kg)}.
				\]
				\item Ta có $\dfrac{n}{4} = 7{,}5$, mà $2<7{,}5<12$ nên nhóm $5$ là nhóm đầu tiên có tần số tích luỹ lớn hơn hoặc bằng $7{,}5$. \\
				Xét nhóm $5$ là nhóm $\left[35;40\right)$ có $s=35$, $h=5$, $n_5=10$ và nhóm $4$ có $cf_4 = 2$.\\
				Khi đó, tứ phân vị thứ nhất là
				\[
				Q_1 = 35 + \dfrac{7{,}5 - 2}{10} \cdot 5 \approx 37{,}8\text{ (kg)}.
				\]
				\item Ta có $\dfrac{3n}{4} = 22{,}5$, mà $12 < 22{,}5 <29$ nên nhóm $6$ là nhóm đầu tiên có tần số tích luỹ lớn hơn hoặc bằng $22{,}5$. \\
				Xét nhóm $6$ là nhóm $\left[40;45\right)$ có $t=40$, $l=5$, $n_6=17$ và nhóm $5$ có $cf_5 = 12$.\\
				Khi đó, tứ phân vị thứ ba là
				\[
				Q_3 = 40 + \dfrac{22{,}5 - 12}{17} \cdot 5 \approx 43{,}1\text{ (kg)}.
				\]
			\end{itemize}
			Vậy $40\text{ (kg)}$ là trung bình mẫu, $40{,}9\text{ (kg)}$ là trung vị của mẫu và $Q_1 = 37{,}8\text{ (kg)}$, $Q_2 = 40{,}9\text{ (kg)}$, $Q_3 = 43{,}1\text{ (kg)}$ là tứ phân vị của mẫu.
			\item Ta thấy: Nhóm $6$ ứng với nửa khoảng $\left[40;45\right)$ là nhóm có tần số lớn nhất với $u=40$, $g=5$, $n_6 = 17$. Nhóm $5$ có tần số $n_5 = 10$, nhóm $7$ có tần số $n_7 = 0$.\\
			Khi đó, mốt của mẫu số liệu là 
			\[
			M_o = 40 + \left( \dfrac{17 - 10}{2\cdot 17 - 10 - 0} \right) \cdot 5 \approx 41{,}5\text{ (kg)}.
			\]
		\end{enumerate}
	}
\end{bt}

\begin{bt}%[1D5H2-3]
% nguon: Archive/2024/TOAN_11_XE/TOAN_11_XE/data4/11C3/KNTT-3.2.tex #925
%[Tex hóa SGK CD, TVN-223]
	\immini
	{
		Bảng bên cho ta bảng tần số ghép nhóm số liệu thống kê chiều cao của $40$ mẫu cây ở một vườn thực vật (đơn vị: centimét).
		\begin{enumerate}
			\item Xác định số trung bình cộng, trung vị, tứ phân vị của mẫu số liệu ghép nhóm trên.
			\item Mốt của mẫu số liệu ghép nhóm trên là bao nhiêu?
		\end{enumerate}
	}
	{
		\begin{tabular}{|c|c|c|}
			\hline
			\textbf{Nhóm} & \textbf{Tần số} & \textbf{Tần số tích luỹ}\\ 
			\hline
			$\left[30;40\right)$ & $4$ & $4$\\
			$\left[40;50\right)$ & $10$ & $14$\\
			$\left[50;60\right)$ & $14$ & $28$\\
			$\left[60;70\right)$ & $6$ & $34$\\
			$\left[70;80\right)$ & $4$ & $38$\\
			$\left[80;90\right)$ & $2$ & $40$\\
			\hline
			& $n = 40$ &\\
			\hline
		\end{tabular}
	}
	\loigiai{
		\begin{enumerate}
			\item 
			\begin{itemize}
				\item Trung bình cộng của mẫu số liệu trên là
				\[
				\overline{x} = \dfrac{35 \cdot 4 + 45 \cdot 10 + 55 \cdot 14 + 65 \cdot 6 + 75 \cdot 4 + 85 \cdot 2}{40} = 	55{,}5\text{ (cm)}.
				\]
				\item Ta có $\dfrac{n}{2} = 20$, mà $14<20<28$ nên nhóm $3$ là nhóm đầu tiên có tần số tích luỹ lớn hơn hoặc bằng $20$. \\
				Xét nhóm $3$ là nhóm $\left[50;60\right)$ có $r=50$, $d=10$, $n_3=14$ và nhóm $2$ có $cf_2 = 14$.\\
				Khi đó, tứ phân vị thứ hai (cũng là trung vị) là
				\[
				Q_2 = M_e = 50 + \dfrac{20 - 14}{14} \cdot 10 = 54{,}3\text{ (cm)}.
				\]
				\item Ta có $\dfrac{n}{4} = 10$, mà $4 < 10 < 14$ nên nhóm $2$ là nhóm đầu tiên có tần số tích luỹ lớn hơn hoặc bằng $10$. \\
				Xét nhóm $2$ là nhóm $\left[40;50\right)$ có $s=40$, $h=10$, $n_2=10$ và nhóm $1$ có $cf_1 = 4$.\\
				Khi đó, tứ phân vị thứ nhất là
				\[
				Q_1 = 40 + \dfrac{10 - 4}{10} \cdot 10 = 46\text{ (cm)}.
				\]
				\item Ta có $\dfrac{3n}{4} = 30$, mà $28 < 30 < 34$ nên nhóm $4$ là nhóm đầu tiên có tần số tích luỹ lớn hơn hoặc bằng $30$. \\
				Xét nhóm $4$ là nhóm $\left[60;70\right)$ có $t = 60$, $l = 10$, $n_4 = 6$ và nhóm $3$ có $cf_3 = 28$.\\
				Khi đó, tứ phân vị thứ ba là
				\[
				Q_3 = 60 + \dfrac{30 - 28}{6} \cdot 10 \approx 63{,}3\text{ (cm)}.
				\]
			\end{itemize}
			Vậy $55{,}5\text{ (cm)}$ là trung bình mẫu, $54{,}3\text{ (cm)}$ là trung vị của mẫu và $Q_1 = 46\text{ (cm)}$, $Q_2 = 54{,}3\text{ (cm)}$, $Q_3 = 63{,}3\text{ (cm)}$ là tứ phân vị của mẫu.
			\item Ta thấy: Nhóm $3$ ứng với nửa khoảng $\left[50;60\right)$ là nhóm có tần số lớn nhất với $u=50$, $g=10$, $n_3 = 14$. Nhóm $2$ có tần số $n_2 = 10$, nhóm $4$ có tần số $n_4 = 6$.\\
			Khi đó, mốt của mẫu số liệu là 
			\[
			M_o = 50 + \left( \dfrac{14 - 10}{2\cdot 14 - 10 - 6} \right) \cdot 10 \approx 53{,}3\text{ (cm)}.
			\]
		\end{enumerate}
	}
\end{bt}

\begin{bt}%[1D9N1-2]
% nguon: Archive/2024/TOAN_11_XE/TOAN_11_XE/data4/11C8/KNTT-8.1.tex #1221
%[TeX hóa Sách 11 CD]%[Hoàng Khắc Ngân]
	Trong hộp kín có $10$ quả bóng màu xanh và $8$ quả bóng màu đỏ, các quả bóng có kích thước và khối lượng giống nhau. Lấy ngẫu nhiên đồng thời $2$ quả bóng. Xét các biến cố:\\
	A: \lq\lq Hai quả bóng lấy ra có màu xanh\rq\rq;\\
	B: \lq\lq Hai quả bóng lấy ra có màu đỏ\rq\rq.\\
	Chọn phát biểu đúng trong những phát biểu sau đây?
	\begin{enumerate}
		\item Biến cố hợp của hai biến cố $A$ và $B$ là \lq\lq Hai quả bóng lấy ra có cùng màu đỏ hoặc cùng màu xanh\rq\rq;
		\item Biến cố hợp của hai biến cố $A$ và $B$ là \lq\lq Hai quả bóng lấy ra có màu khác nhau\rq\rq;
		\item Biến cố hợp của hai biến cố $A$ và $B$ là \lq\lq Hai quả bóng lấy ra có cùng màu\rq\rq.
	\end{enumerate}
	\loigiai{\begin{multicols}{3}
		\begin{enumerate}
			\item Phát biểu đúng;
			\item Phát biểu sai;
			\item Phát biểu đúng.
		\end{enumerate}
		\end{multicols}	
	}
\end{bt}

\begin{bt}%[1D9N2-2]
% nguon: Archive/2024/TOAN_11_XE/TOAN_11_XE/data4/11C8/KNTT-8.1.tex #1222
%[TeX hóa Sách 11 CD]%[Hoàng Khắc Ngân]
	Một hộp có $52$ chiếc thẻ cùng loại, mỗi thẻ được ghi một trong các số $1,2,3, \ldots, 52$; hai thẻ khác nhau thì ghi hai số khác nhau. Rút ngẫu nhiên $1$ chiếc thẻ trong hộp. Xét biến cố $A$: \lq\lq Số xuất hiện trên thẻ được rút ra là số chia hết cho 3\rq\rq \, và biến cố $B$: \lq\lq Số xuất hiện trên thẻ được rút ra là số lẻ\rq\rq. Viết các tập con của không gian mẫu tương ứng với các biến cố $A, B, A \cap B$.
	\loigiai{
		Ta có $A=\{3 ; 6 ; 9 ; 12 ; 15 ; \ldots ; 48 ; 51\}$;\\
		$B=\{4 ; 8 ; 12 ; 16 ; 20 ; \ldots ; 48 ; 52\}$;\\
		$A \cap B=\{12 ; 24 ; 36 ; 48\}$.
	}
\end{bt}

\begin{bt}%[1D9N2-2]
% nguon: Archive/2024/TOAN_11_XE/TOAN_11_XE/data4/11C8/KNTT-8.1.tex #1223
%[TeX hóa Sách 11 CD]%[Hoàng Khắc Ngân]
	Một hộp có $3$ quả bóng màu xanh, $4$ quả bóng màu đỏ; các quả bóng có kích thước và khối lượng như nhau. Lấy bóng ngẫu nhiên hai lần liên tiếp, trong đó mỗi lần lấy ngẫu nhiên một quả bóng trong hộp, ghi lại màu của quả bóng lấy ra và bỏ lại quả bóng đó vào hộp.\\
	Xét các biến cố:\\
	$A$: \lq\lq Quả bóng màu xanh được lấy ra ở lần thứ nhất\rq\rq;\\
	$B$: \lq\lq Quả bóng màu đỏ được lấy ra ở lần thứ hai\rq\rq.\\
	Hỏi
	\begin{enumerate}[a)]
		\item Hai biến cố $A$ và $B$ có độc lập không? Vì sao?
		\item Hai biến cố $A$ và $B$ có xung khắc không? Vì sao?
	\end{enumerate}
	\loigiai{
		\begin{enumerate}[a)]
			\item Trước hết, biến cố $B$ xảy ra sau biến cố $A$ nên việc xảy ra hay không xảy ra của biến cố $B$ không làm ảnh hưởng đến xác suất xảy ra của biến cố $A$.\\
			Mặt khác, ta có xác suất của biến cố $B$ khi biến cố $A$ xảy ra bằng $\dfrac{4}{7}$; xác suất của biến cố $B$ khi biến cố $A$ không xảy ra cũng bằng $\dfrac{4}{7}$.\\
			Do đó việc xảy ra hay không xảy ra của biến cố $A$ không làm ảnh hưởng đến xác suất xảy ra của biến cố $B$.\\
			Vậy hai biến cố $A$ và $B$ là độc lập.
			\item Ta thấy kết quả (xanh; đỏ) là kết quả thuận lợi cho cả hai biến cố $A$ và $B$. Vì thế $A$ và $B$ không là hai biến cố xung khắc.
	\end{enumerate}}		
\end{bt}

\begin{bt}%[1D9N1-2]
% nguon: Archive/2024/TOAN_11_XE/TOAN_11_XE/data4/11C8/KNTT-8.1.tex #1224
%[TeX hóa Sách 11 CD]%[Hoàng Khắc Ngân]
	Tung một đồng xu cân đối và đồng chất hai lần liên tiếp. Xét các biến cố:\\
	$A$: \lq\lq Lần thứ nhất xuất hiện mặt ngửa\rq\rq;\\
	$B$: \lq\lq Lần thứ hai xuất hiện mặt ngửa\rq\rq;\\
	$C$: \lq\lq Cả hai lần đều xuất hiện mặt ngửa\rq\rq;\\
	$D$: \lq\lq Có ít nhất một lần xuất hiện mặt ngử\rq\rq.\\
	Trong hai biến cố $C, D$, biến cố nào là biến cố hợp của hai biến cố $A, B$ ? Biến cố nào là biến cố giao của hai biến cố $A, B$?
	\loigiai{
		Biến cố $C$ là biến cố giao của hai biến cố $A, B$.\\
		Biến cố $D$ là biến cố hợp của hai biến cố $A, B$.
	}
\end{bt}

\begin{bt}%[1D9N2-2]
% nguon: Archive/2024/TOAN_11_XE/TOAN_11_XE/data4/11C8/KNTT-8.1.tex #1225
%[TeX hóa Sách 11 CD]%[Hoàng Khắc Ngân]
	Gieo ngẫu nhiên một xúc xắc cân đối và đồng chất hai lần liên tiếp. Xét các biến cố:\\
	$A$: \lq\lq Số chấm xuất hiện ở lần gieo thứ nhất lớn hơn 4\rq\rq ;\\
	$B$: \lq\lq Số chấm xuất hiện ở lần gieo thứ hai nhỏ hơn 4\rq\rq ;\\
	$C$: \lq\lq Số chấm xuất hiện ở lần gieo thứ nhất nhỏ hơn 4\rq\rq.\\
	Trong các biến cố trên, hãy:
	\begin{multicols}{2}
		\begin{enumerate}
			\item Tìm cặp biến cố xung khắc;
			\item Tìm cặp biến cố độc lập.
		\end{enumerate}
	\end{multicols}
	\loigiai{
		\begin{enumerate}
			\item Cặp biến cố xung khắc là $A, C$.
			\item Cặp biến cố độc lập là $A, B$ và $B, C$.
		\end{enumerate}
	}
\end{bt}

\begin{bt}%[1D9H2-4]
% nguon: Archive/2024/TOAN_11_XE/TOAN_11_XE/data4/11C8/KNTT-8.2.tex #1228
%[TeX hóa Sách 11 CD]%[Hoàng Khắc Ngân]
	Chọn ngẫu nhiên một số nguyên dương có hai chữ số. Xét biến cố $A$: \lq\lq Số được viết ra là số chia hết cho $8$\rq\rq \, và biến cố $B$:  \lq\lq Số được viết ra là số chia hết cho $9$\rq\rq. Tính $\mathrm{P}(A \cup B)$.
	\loigiai{Trong $90$ số có hai chữ số thì
		\begin{enumerate}[+)]
			\item có $11$ số chia hết cho $8$,
			\item có $10$ số chia hết cho $9$,
			\item có $1$ số chia hết cho cả $8$ và $9$.
		\end{enumerate} 
		Vì thế, ta có $\mathrm{P}(A)=\dfrac{11}{90}, \mathrm{P}(B)=\dfrac{10}{90}, \mathrm{P}(A \cap B)=\dfrac{1}{90}$.\\
		Vậy $\mathrm{P}(A \cup B)=\mathrm{P}(A)+\mathrm{P}(B) - \mathrm{P}(A \cap B) =\dfrac{11}{90}+\dfrac{10}{90}-\dfrac{1}{90}=\dfrac{2}{9}$.}
\end{bt}

\begin{bt}%[1D9H2-4]
% nguon: Archive/2024/TOAN_11_XE/TOAN_11_XE/data4/11C8/KNTT-8.2.tex #1229
%[TeX hóa Sách 11 CD]%[Hoàng Khắc Ngân]
	Một hộp có $12$ chiếc thẻ cùng loại, mỗi thẻ được ghi một trong các số $1,2,3, \ldots, 12$; hai thẻ khác nhau thì ghi hai số khác nhau. Rút ngẫu nhiên $1$ chiếc thẻ trong hộp. Xét biến cố $A$: \lq\lq Số xuất hiện trên thẻ được rút ra là số chia hết cho $3$\rq\rq \, và biến cố $B$: \lq\lq Số xuất hiện trên thẻ được rút ra là số chia hết cho 5\rq\rq. Tính $\mathrm{P}(A \cup B)$.
	\loigiai{
		Không gian mẫu của phép thử trên có $12$ phần tử, tức là $n(\Omega)=12$.\\
		Số các kết quả thuận lợi cho các biến cố $A, B$ lần lượt là $n(A)=4, n(B)=2$.\\
		Suy ra
		$$\mathrm{P}(A)=\dfrac{n(A)}{n(\Omega)}=\dfrac{4}{12}=\dfrac{1}{3}, \mathrm{P}(B)=\dfrac{n(B)}{n(\Omega)}=\dfrac{2}{12}=\dfrac{1}{6}.$$
		Trong các số $1,2,3, \ldots, 12$, không có số nào chia hết cho cả $3$ và $5$.\\
		Vì thế $A, B$ là hai biến cố xung khắc. Suy ra
		$\mathrm{P}(A \cup B)=\mathrm{P}(A)+\mathrm{P}(B)=\dfrac{1}{3}+\dfrac{1}{6}=\dfrac{1}{2}$.}
\end{bt}

\begin{bt}%[1D9N2-3]
% nguon: Archive/2024/TOAN_11_XE/TOAN_11_XE/data4/11C8/KNTT-8.2.tex #1230
%[TeX hóa Sách 11 CD]%[Hoàng Khắc Ngân]
	Chọn ngẫu nhiên một số tự nhiên có hai chữ số. Tính xác suất của biến cố $M$: \lq\lq Số tự nhiên có hai chữ số được viết ra chia hết cho $11$ hoặc chia hết cho $12$ \rq\rq.
	\loigiai{
		Số các số tự nhiên có hai chữ số là $90$ số. Do đó số cách chọn ngẫu nhiên một số tự nhiên có hai chữ số là $90$ nên $ n(\Omega)=90$.\\
		Số các số tự nhiên có hai chữ số chia hết cho $11$ hoặc chia hết cho $12$ là $18$ số (gồm có $11$, $22$, $33$, $44$, $55$, $66$, $77$, $88$, $99$, $12$, $24$, $36$, $48$, $50$, $62$, $74$, $86$, $98$), nên $n(M)=18$.\\
		Vậy xác suất của biến cố $M$ là $\mathrm{P}(M)=\dfrac{n(M)}{n(\Omega)}=\dfrac{18}{90}=\dfrac{1}{5}$.
	}
\end{bt}

\begin{bt}%[1D9H1-3]
% nguon: Archive/2024/TOAN_11_XE/TOAN_11_XE/data4/11C8/KNTT-8.3.tex #1232
%[TeX hóa Sách 11 CD]%[Hoàng Khắc Ngân]
	Hai bạn Hạnh và Hà cùng chơi trò chơi bắn cung một cách độc lập. Mỗi bạn chỉ bắn một lần. Xác suất để bạn Hạnh và bạn Hà bắn trúng bia lần lượt là $0{,}6$ và $0{,}7$ trong lần bắn của mình. Tính xác suất của biến cố $C$: \lq\lq Bạn Hạnh và bạn Hà đều bắn trúng bia \rq\rq.
	\loigiai{
		Xét biến cố $A$: \lq\lq Bạn Hạnh bắn trúng bia\rq\rq, ta có $\mathrm{P}(A)=0{,}6$.\\
		Xét biến cố $B$: \lq\lq Bạn Hà bắn trúng bia\rq\rq, ta có $\mathrm{P}(B)=0{,}7$.\\
		Ta thấy $A, B$ là hai biến cố độc lập và $C=A \cap B$. suy ra
		$$\mathrm{P}(C)=\mathrm{P}(A) \cdot \mathrm{P}(B)=0{,}6 \cdot 0{,}7=0{,}42.$$}
\end{bt}

\begin{bt}%[1D9H1-3]
% nguon: Archive/2024/TOAN_11_XE/TOAN_11_XE/data4/11C8/KNTT-8.3.tex #1233
%[TeX hóa Sách 11 CD]%[Hoàng Khắc Ngân]
	Hai bạn Trung và Dũng của lớp $11A$ tham gia giải bóng bàn đơn nam do nhà trường tổ chức. Hai bạn đó không cùng thuộc một bảng đấu loại và mỗi bảng đấu loại chỉ chọn một người vào vòng chung kết. Xác suất lọt qua vòng loại để vào vòng chung kết của Trung và Dũng lần lượt là $0{,}8$ và $0{,}6$. Tính xác suất của các biến cố sau
	\begin{enumerate}[a)]
		\item $A$: \lq\lq Cả hai bạn lọt vào vòng chung kết\rq\rq;
		\item $B$: \lq\lq Có ít nhất một bạn lọt vào vòng chung kết \rq\rq;
		\item $C$: \lq\lq Chỉ có bạn Trung lọt vào vòng chung kết\rq\rq.		
	\end{enumerate}
	
	\loigiai{
		Xét các biến cố $E$: \lq\lq Bạn Trung lọt vào vòng chung kết\rq\rq \, và $G$: \lq\lq Bạn Dũng lọt vào vòng chung kết\rq\rq.\\
		Từ giả thiết, ta suy ra $E, G$ là hai biến cố độc lập và $\mathrm{P}(E)=0{,}8; \mathrm{P}(G)=0{,}6$.\\
		\begin{enumerate}[a)]
			\item Do $A=E \cap G$ nên $\mathrm{P}(A)=\mathrm{P}(E) \cdot \mathrm{P}(G)=0{,}8 \cdot 0{,}6=0{,}48$.
			\item Ta thấy $B=E \cup G$, suy ra
			$$	\mathrm{P}(B)=\mathrm{P}(E \cup G)=\mathrm{P}(E)+\mathrm{P}(G)-\mathrm{P}(E \cap G)=0{,}8+0{,}6-0{,}48=0{,}92.$$
			\item Xét biến cố đối $\bar{G}$ của biến cố $G$.\\
			Ta thấy $\mathrm{P}(\bar{G})=1-\mathrm{P}(G)=1-0{,}6=0{,}4$ và $E, \bar{G}$ là hai biến cố độc lập.\\
			Vì $C=E \cap \bar{G}$ nên $\mathrm{P}(C)=\mathrm{P}(E) \cdot \mathrm{P}(\bar{G})=0{,}8 \cdot 0{,}4=0{,}32$.
		\end{enumerate}
	}
\end{bt}

\begin{bt}%[1D9N2-2]
% nguon: Archive/2024/TOAN_11_XE/TOAN_11_XE/data4/11C8/KNTT-8.3.tex #1235
%[TeX hóa Sách 11 CD]%[Hoàng Khắc Ngân]
	Gieo ngẫu nhiên một xúc xắc cân đối và đồng chất hai lần liên tiếp. Xét các biến cố:\\
	$A$: \lq\lq Số chấm xuất hiện ở lần gieo thứ nhất lớn hơn 4\rq\rq ;\\
	$B$: \lq\lq Số chấm xuất hiện ở lần gieo thứ hai nhỏ hơn 4\rq\rq ;\\
	$C$: \lq\lq Số chấm xuất hiện ở lần gieo thứ nhất nhỏ hơn 4\rq\rq.\\
	Trong các biến cố trên, hãy:
	\begin{enumerate}[a)]
		\item Tìm cặp biến cố xung khắc;
		\item Tìm cặp biến cố độc lập.
	\end{enumerate}
	\loigiai{
		\begin{enumerate}[a)]
			\item Cặp biến cố xung khắc là $A, C$.
			\item Cặp biến cố độc lập là $A, B$ và $B, C$.
		\end{enumerate}
	}
\end{bt}

\begin{bt}%[1D9V1-3]
% nguon: Archive/2024/TOAN_11_XE/TOAN_11_XE/data4/11C8/KNTT-8.3.tex #1236
%[TeX hóa Sách 11 CD]%[Hoàng Khắc Ngân]
	Hai bạn Việt và Nam cùng tham gia một kì thi trắc nghiệm môn Toán và môn Tiếng Anh một cách độc lập nhau. Đề thi của mỗi môn gồm $6$ mã đề khác nhau và các môn khác nhau thì mã đề cũng khác nhau. Đề thi được sắp xếp và phát cho học sinh một cách ngẫu nhiên. Tính xác suất để hai bạn Việt và Nam có chung đúng một mã đề thi trong kì thi đó.
	\loigiai{
		Việt có số cách nhận mã đề hai môn là $6 \cdot 6 = 36$ cách.\\
		Nam có số cách nhận mã đề hai môn là $6 \cdot 6 = 36$ cách.\\
		Số phần tử của không gian mẫu là $36 \cdot 36 = 1296$.\\
		Gọi $H$ là biến cố \lq\lq Việt và Nam có chung đúng một mã đề thi\rq\rq.\\
		Có các trường hợp sau:
		\begin{enumerate}
			\item Có cùng mã đề môn Toán: Việt có $36$ cách nhận mã đề hai môn, Nam có $5$ cách nhận mã đề. Do đó có $36 \cdot 5 = 180$ cách.
			\item Tương tự có cùng môn Tiếng Anh có $180$ cách.
		\end{enumerate}
		Suy ra $n(H)=360$ cách.\\
		Vậy xác suất của biến cố $H$ là $\mathrm{P}(H)=\dfrac{n(H)}{n(\Omega)}=\dfrac{360}{1296}=\dfrac{5}{18}$.
	}
\end{bt}

\begin{bt}%[1D9V2-3]
% nguon: Archive/2024/TOAN_11_XE/TOAN_11_XE/data4/11C8/KNTT-8.3.tex #1237
%[TeX hóa Sách 11 CD]%[Hoàng Khắc Ngân]
	Trong một chiếc hộp có $20$ viên bi có cùng kích thước và khối lượng, trong đó có $9$ viên bi màu đỏ, $6$ viên bi màu xanh và $5$ viên bi màu vàng. Lấy ngẫu nhiên đồng thời $3$ viên bi. Tîm xác suất để $3$ viên bi lấy ra có đúng hai màu.
	\loigiai{
		Mỗi cách lấy $3$ viên bi trong $20$ viên bi là một tổ hợp chập $3$ của $20$ phần tử.\\
		Do đó, không gian mẫu $\Omega$ gồm các tổ hợp chập $3$ của $20$ phần tử và $ n(\Omega)=\mathrm{C}_{20}^3=1140$.\\
		Lấy $3$ viên bi có đủ ba màu là $9 \cdot 6 \cdot 5 = 270$ cách.\\
		Lấy $3$ viên bi có đúng một màu là $\mathrm{C}_9^3 + \mathrm{C}_6^3 + \mathrm{C}_5^3 = 114$ cách.\\
		Do đó lấy $3$ viên bi có đúng hai màu là $\mathrm{C}_{20}^3 - (270 + 114) = 756$ cách.\\
		Vậy xác suất của biến cố để $3$ viên bi lấy ra có đúng hai màu là $\mathrm{P}(H)=\dfrac{n(H)}{n(\Omega)}=\dfrac{756}{1140}=\dfrac{63}{95}$.
	}
\end{bt}

\begin{bt}%[1D5H1-3]
% nguon: Archive/2024/Update-SP-BG10-11-3/SP-3/BG11-3/data_1C5/ChuongV-Bai1-Dang2-SoTrungBinhCong-Diamon.tex #1500
%[KNTT]
Tìm cân nặng trung bình của học sinh lớp $11D$ cho trong Bảng $3.5$.
\begin{center}
\begin{tabular}{|c|c|c|c|c|c|c|}
\hline
Cân nặng& $\left[40{,}5;45{,}5 \right)$ & $\left[45{,}5;50{,}5 \right)$ & $\left[50{,}5;55{,}5 \right)$ & $\left[55{,}5;60{,}5 \right)$ & $\left[60{,}5;65{,}5 \right)$ & $\left[65{,}5;70{,}5 \right)$ \\
\hline
Số học sinh&$10$& $7$ & $16$ &$4$ & $2$ & $3$ \\
\hline
\end{tabular}\\
\vspace*{0.5cm}
Bảng $3.5$. Cân nặng của học sinh lớp $11D$.
\end{center}
\loigiai{
Trong mỗi khoảng cân nặng, giá trị đại diện là trung bình cộng của hai giá trị đầu mút nên ta có bảng sau
\begin{center}
\begin{tabular}{|c|c|c|c|c|c|c|}
\hline
Cân nặng (kg)& $43$ & $48$ & $53$ & $58$ & $63$ & $68$ \\
\hline
Số học sinh &$10$ & $7$ & $16$ &$4$ & $2$ & $3$ \\
\hline
\end{tabular}
\end{center}
Tổng số học sinh là $n=42$. Cân nặng trung bình của học sinh lớp $11D$ là $$\overline{x}=\dfrac{10\cdot 43+7\cdot 48+16\cdot 53+4\cdot 58+2\cdot 63+3\cdot 68}{42}\approx51{,}81\,\mathrm{(kg)}.$$
}
\end{bt}

\begin{bt}%[1D5H1-3]
% nguon: Archive/2024/Update-SP-BG10-11-3/SP-3/BG11-3/data_1C5/ChuongV-Bai1-Dang2-SoTrungBinhCong-Diamon.tex #1502
%[KNTT]
Tuổi thọ (năm) của 50 bình ắc quy ô tô được cho như sau:
\begin{center}
\begin{tabular}{|c|c|c|c|c|c|c|}
\hline
Tuổi thọ (năm)& $\left[2;2{,}5 \right)$ & $\left[2{,}5;3 \right)$ & $\left[3;3{,}5 \right)$&$\left[3{,}5;4 \right)$&$\left[4;4{,}5 \right)$&$\left[4{,}5;5 \right)$ \\
\hline
Tần số &$4$& $9$ & $14$ &$11$ & $7$&$5$ \\
\hline
\end{tabular}
\end{center}
Tính tuổi thọ trung bình của $50$ bình ắc quy ô tô này.
\loigiai{
Ta có bảng sau
\begin{center}
\begin{tabular}{|c|c|c|c|c|c|c|}
\hline
Tuổi thọ (năm)& $2{,}25$ & $2{,}75$ & $3{,}25$&$3{,}75$&$4{,}25$&$4{,}75$ \\
\hline
Tần số &$4$& $9$ & $14$ &$11$ & $7$&$5$\\
\hline
\end{tabular}
\end{center}
Tuổi thọ trung bình của 50 bình ắc quy ô tô này là
$$\overline{x}=\dfrac{2{,}25\cdot 4+2{,}75\cdot 9+3{,}25\cdot 14+3{,}75\cdot 11+4{,}25\cdot 7+4{,}75\cdot 5}{50}=3{,}48 \, \text{(năm)}.$$
}
\end{bt}

\begin{bt}%[1D5H1-3]
% nguon: Archive/2024/Update-SP-BG10-11-3/SP-3/BG11-3/data_1C5/ChuongV-Bai1-Dang2-SoTrungBinhCong-Diamon.tex #1503
%[KNTT]
Phỏng vấn một số học sinh lớp $11$ về thời gian (giờ) ngủ của một buổi tối, thu được bảng số liệu ở bên. So sánh thời gian ngủ trung bình của các bạn học sinh nam và nữ.
\begin{center}
\begin{tabular}{|c|c|c|c|c|c|}
\hline
Thời gian& $[4;5)$&$[5;6)$&$[6;7)$&$[7;8)$&$[8;9)$\\
\hline
Số học sinh nam & $6$&$10$& $13$& $9$& $7$ \\
\hline
Số học sinh nữ& $4$& $8$& $10$& $11$& $8$ \\
\hline
\end{tabular}
\end{center}
\loigiai{
Trong mỗi khoảng thời gian, giá trị đại diện là trung bình cộng của giá trị hai đầu mút nên ta có bảng sau:
\begin{center}
\begin{tabular}{|c|c|c|}
\hline
Thời gian& Số học sinh nam & Số học sinh nữ\\
\hline
$4{,}5$& $6$ & $4$ \\
\hline
$5{,}5$& $10$ & $8$ \\
\hline
$6{,}5$& $13$ & $10$ \\
\hline
$7{,}5$& $9$ & $11$ \\
\hline
$8{,}5$& $7$ & $8$ \\
\hline
\end{tabular}
\end{center}
Tổng số học sinh nam là $n_1=6+10+13+9+7=45$.\\ Thời gian ngủ trung bình của học sinh nam là
$$\overline{x_1}=\dfrac{4{,}5\cdot 6+5{,}5\cdot10+6{,}5\cdot13+7{,}5\cdot9+8{,}5\cdot7}{45}=\dfrac{587}{90}\approx 6{,}52\,\, \text{(giờ)}.$$
Tổng số học sinh nữ là $n_2=4+8+10+11+8=41$. Thời gian ngủ trung bình của học sinh nữ là
$$\overline{x_2}=\dfrac{4,5\cdot4+5,5\cdot8+6,5\cdot10+7,5\cdot11+8,5\cdot8}{41}=\dfrac{555}{82}\approx 6{,}77 \,\,\text{(giờ)}.$$
Vì $\overline{x_2}>\overline{x_1}$ nên thời gian ngủ trung bình của các bạn học sinh nữ lớn hơn thời gian ngủ trung bình của các bạn nam.
}
\end{bt}

\begin{bt}%[1D5H1-2]
% nguon: Archive/2024/Update-SP-BG10-11-3/SP-3/BG11-3/data_1C5/ChuongV-Bai1-Dang1-MauSoLieuGhepNhom-Diamon.tex #1529
%[ dự án CTST-Bài giảng 11 - new]%%
Một thư viện thống kê số người đến đọc sách vào buổi tối trong 30 ngày của tháng vừa qua như sau
\begin{center}
$\begin{array}{llllllllllllllll}
&85 & 81 & 65 & 58 & 47 & 30 & 51 & 92 & 85 & 42 & 55 & 37 & 31 & 82 & 63 \\
& 33 & 44 & 93 & 77 & 57 & 44 & 74 & 63 & 67 & 46 & 73 & 52 & 53 & 47 & 35
\end{array}$
\end{center}
Lập bảng tần số ghép nhóm có tám nhóm ứng với tám nửa khoảng sau
\begin{center}
$[25 ; 34)$; $[34 ; 43)$; $[43 ; 52)$; $[52 ; 61)$; $[61 ; 70)$; $[70 ; 79)$; $[79 ; 88)$; $[88 ; 97)$.
\end{center}
\loigiai{
Dựa vào mẫu số liệu và gợi ý theo các nhóm, chúng ta đếm tấn số mỗi nhóm và ghi lại như bảng sau: (Cẩn thận và nên đánh dấu các giá trị đã đếm, tránh nhầm lẫn)
\begin{center}
\begin{tabular}{|c|c|c|c|c|c|c|c|c|}
\hline
Nhóm &$[25 ; 34)$  &$[34 ; 43)$&$[43 ; 52)$&$[52 ; 61)$&$[61 ; 70)$&$[70 ; 79)$ &$[79 ; 88)$& $[88 ; 97)$\\
\hline
Tần số &$3$&$3$  &$5$  &$6$ & $4$ &$3$ &$4$&$2$ \\
\hline
\end{tabular}
\end{center}
}
\end{bt}

\begin{bt}%[1D5N1-2]
% nguon: Archive/SP_KhoLatex/ID6_old/1K3-8.tex #2648

Mẫu số liệu sau cho biết phân bố theo độ tuổi của dân số Việt Nam năm $2019$.
\begin{center}
    \begin{tabular}{|c|c|c|c|}
        \hline
Độ tuổi&Dưới $15$&Từ $15$ đến dưới $65$ tuổi&Từ $65$ tuổi trở lên\\
\hline
Số người&$23371882$&$65420451$&$7416651$\\
\hline
    \end{tabular}
\end{center}
\begin{enumerate}
    \item Mẩu số liệu đã cho có là mẫu số liệu ghép nhóm hay không?
\item Nêu các nhóm và tần số tương ứng. Dân số Việt Nam năm 2019 là bao nhiêu?
\end{enumerate}
\loigiai{
    \begin{enumerate}
        \item Mẫu số liệu đã cho là mẫu số liệu ghép nhóm.
\item 
\begin{itemize}
    \item Có ba nhóm là: Dưới $15$ tuồi, Từ $15$ đến dưới $65$ tuổi, Từ $65$ tuổi trở lên.
\item Có $23371882$ người dưới $15$ tuổi; $65420451$ người từ $15$ đến dưới $65$ tuổi và $7416651$ người từ $65$ tuổi trở lên.
\item  Dân số Việt Nam năm $2019$ là $23371882 + 65420451 + 7416651 = 96208984$ người.

\end{itemize}
\end{enumerate}
}
\end{bt}

\begin{bt}%[1D5N1-2]
% nguon: Archive/SP_KhoLatex/ID6_old/1K3-8.tex #2651

Thời gian ra sân (giờ) của một số cựu cầu thủ ở giải ngoại hạng Anh qua các thời kì được cho như sau:
\begin{center}
    \begin{tabular}{llllllllllllllll}
        $653$&$632$&$609$&$572$&$565$&$535$&$516$&$514$&$508$&$505$&$504$&$504$&$503$&$499$&$496$&$492$
    \end{tabular}
\begin{flushright}
    (Theo :https://www.premierleague.com/)
\end{flushright}
\end{center}
Hãy chuyển mẫu số liệu trên sang dạng ghép nhóm với bảy nhóm có độ dài bằng nhau.
\loigiai{
Giá trị lớn nhất là: $653$.\\
Giá trị bé nhất là: $492$.\\
Khoảng biến thiên là: $653-429=161$.  \\
Để chia thành $7$ nhóm có độ dài bằng nhau, ta lấy điểm đầu mút trái của nhóm đầu tiên là $492$, điểm đầu mút phải của nhóm cuối là $653$ với độ dài mỗi nhóm là $23$.\\
Ta có mẫu số liệu ghép nhóm sau:
\begin{center}
    \begin{tabular}{|c|c|c|c|c|c|c|c|}
    \hline
    Thời gian ra sân (giờ)&$[492;515)$&$[515;538)$&$[538;561)$&$[561;584)$&$[584;607)$&$[607;630)$&$[630;653)$\\
    \hline
    Số cầu thủ&$9$&$2$&$0$&$2$&$0$&$1$&$2$\\
    \hline
    \end{tabular}
\end{center}
}
\end{bt}

\begin{bt}%[1D5H1-2]
% nguon: Archive/SP_KhoLatex/ID6_old/1K3-8.tex #2652

Cho các số liệu thống kê ghi trong bảng sau: 
Thành tích chạy $50$cm của học sinh lóp $10$A  ở trường Trung học phổ thông C (đơn vị : giây)
\begin{center}
    \begin{tabular}{|lllllll|}
    \hline
    $6{,}3$&$6{,}2$&$6{,}5$&$6{,}8$&$6{,}9$&$8{,}2$&$8{,}6$\\
        $6{,}6$&$6{,}7$&$7{,}0$&$7{,}1$&$7{,}2$&$8{,}3$&$8{,}5$\\
            $7{,}4$&$7{,}3$&$7{,}2$&$7{,}1$&$7{,}0$&$8{,}4$&$8{,}1$\\
                $7{,}1$&$7{,}3$&$7{,}5$&$7{,}5$&$7{,}6$&$8{,}7$&\\
                $7{,}6$&$7{,}7$&$7{,}8$&$7{,}5$&$7{,}7$&$7{,}8$&\\
                \hline
    \end{tabular}
\end{center}
\begin{enumerate}
\item Lập bảng phân bố tần số ghép nhóm và bảng phân bố tần suất ghép nhóm, với các nhóm $[6{,}0;6{,}5)$, $[6{,}5;7{,}0)$, $[7{,}0;7{,}5)$,$[7{,}5;8{,}0)$,$[8{,}0;8{,}5)$,$[8{,}5;9{,}0)$.
\item Trong lớp $10$A, số học sinh chạy $50$m  hết từ $7$ giây đến dưới $8{,}5$ giây chiếm bao nhiêu phần trăm ?
\end{enumerate}
\loigiai{
    \begin{enumerate}
        \item Từ các số liệu thống kê đã cho, ta xác định được:\\
Tần số của các lớp: $$n_1=2; n_2=5; n_3=10; n_4=9; n_5=4; n_6=3.$$

Tần suất của các lớp: 
\[f_1\approx 6{,}06\%; f_2\approx 15{,}15\%; f_3\approx 30{,}30\%; f_4\approx 27{,}27\%; f_5\approx 12{,}12\%; f_6\approx 9{,}10\% \]
Từ đó ta có bảng phân bố tần số ghép nhóm: 
Thành tích chay $50$ m của học sinh lớp $10$ A ở truơng Trung học phổ thông C
\begin{center}
    \begin{tabular}{|c|c|}
    \hline
    Lớp thời gian chạy (giây)&Tần suất ($\%$)\\
    \hline
    $[6{,}0;6{,}5)$&$6{,}06$\\
     $[6{,}5;7{,}0)$&$15{,}15$\\ $[7{,}0;7{,}5)$&$30{,}30$\\
     $[7{,}5;8{,}0)$&$27{,}27$\\
     $[8{,}0;8{,}5)$&$12{,}12$\\
     $[8{,}5;9{,}0)$&$9{,}10$\\
     \hline
     Cộng&$100\%$\\
     \hline
    \end{tabular}
\end{center}

b) $30{,}30\%+27{,}27\%+12{,}12\%= 69{,}69\%$\\  .
Trả lời : $69{,}69\%$.
\end{enumerate}
}
\end{bt}

\begin{bt}%[1D5H2-3]
% nguon: Archive/SP_KhoLatex/ID6_old/1K3-9.tex #2655

    Thời gian (phút) truy cập Internet mỗi buổi tối của một số học sinh được cho trong bảng sau:
    \begin{center}
        \begin{tabular}{|c|c|c|c|c|c|c|} 
            \hline
            Thời gian (giờ) & $[9{,}5; 12{,}5)$ & $[12{,}5; 15{,}5)$ & $[15{,}5; 18{,}5)$ & $[18{,}5; 21{,}5)$ & $[21{,}5; 24{,}5]$ \\
            \hline
            Số học sinh & $3$ & $12$ & $15$ & $24$ & $2$ \\
            \hline
        \end{tabular}
    \end{center}
    Tìm tứ phân vị thứ nhất $Q_1$ và tứ phân vị thứ ba $Q_3$ của mẫu số liệu ghép nhóm cho đa cho.
    \loigiai{
    Cỡ mẫu là $n = 56$.\\
    Tứ phân vị thứ nhất $Q_1$ là $\dfrac{x_{14} + x_{15}}{2}$.\\
    Do $x_{14}$, $x_{15}$ đều thuộc nhóm $[12{,}5; 15{,}5)$ nên nhóm này chứa $Q_1$.\\
    Do đó $p = 2$, $a_2 = 12{,}5$; $m_2 = 12$; $m_1 = 3$; $a_3 - a_2 = 3$ và ta có 
    $$Q_1 = 12{,}5 + \dfrac{\dfrac{56}{4} - 3}{12} \cdot 3 = 15{,}25.$$
    Tứ phân vị thứ ba $Q_3$ là $\dfrac{x_{42} + x_{43}}{2}$.\\
    Do $x_{42}$, $x_{43}$ đều thuộc nhóm $[18{,}5; 21{,}5)$ nên nhóm này chứa $Q_3$.\\
    Do đó $p = 4$, $a_2 = 18{,}5$; $m_4 = 24$; $m_1 + m_2 + m_3 = 3 + 12 + 15 = 30$; $a_5 - a_4 = 3$ và ta có $$Q_3 = 18{,}5 + \dfrac{\dfrac{3 \cdot 56}{4} - 30}{24} \cdot 3 = 20.$$
    }
\end{bt}

\begin{bt}%[1D5H1-4]
% nguon: Archive/SP_KhoLatex/ID6_old/1K3-9.tex #2658

    Tuổi thọ (năm) của $50$ bình ắc quy ô tô được cho như sau:
    \begin{center}
        \begin{tabular}{|c|c|c|c|c|c|c|} 
            \hline
            Tuổi thọ (năm) & $[2; 2{,}5)$ & $[2{,}5; 3)$ & $[3; 3{,}5)$ & $[3{,}5; 4)$ & $[4; 4{,}5)$ & $[4{,}5; 5]$\\
            \hline
            Tần số & $4$ & $9$ & $14$ & $11$ & $7$ & $5$\\
            \hline
        \end{tabular}
    \end{center}
    \begin{enumerate}
        \item Xác định mốt và giải thích ý nghĩa.
        \item Tính tuổi thọ trung bình của $50$ bình ắc quy ô tô này.
    \end{enumerate}
    \loigiai{
    \begin{enumerate}
        \item $14$ là tần số lớn nhất nên mốt thuộc nhóm $[3; 3{,}5)$, ta có $j = 3$, $a_3 = 3$, $m_3 = 14$, $_2 = 9$, $m_4 = 11$, $h = 0{,}5$.\\
        Do đó $M_0 = 3 + \dfrac{14 - 9}{(14 - 9) + (14 - 11)} \cdot 0{,}5  = 3{,}31$.
        \item Bảng giá trị đại diện
        \begin{center}
            \begin{tabular}{|c|c|c|c|c|c|c|} 
                \hline
                Tuổi thọ (năm) & $[2; 2{,}5)$ & $[2{,}5; 3)$ & $[3; 3{,}5)$ & $[3{,}5; 4)$ & $[4; 4{,}5)$ & $[4{,}5; 5]$\\
                \hline
                Tần số & $4$ & $9$ & $14$ & $11$ & $7$ & $5$\\
                \hline
                Tần số & $2{,}25$ & $2{,}75$ & $3{,}25$ & $3{,}75$ & $4{,}25$ & $4{,}75$\\
                \hline
            \end{tabular}
        \end{center}
        Tuổi thọ trung bình của $50$ bình ắc quy ô tô là 
        $$\overline{x} = \dfrac{4 \cdot 2{,}25 + 9 \cdot 2{,}75 + 14 \cdot 3{,}25 + 11 \cdot 3{,}75 + 7 \cdot 4{,}25 + 5 \cdot 4{,}75 + }{50} = 30{,}48.$$
    \end{enumerate}
    }
\end{bt}

\begin{bt}%[1D5H1-3]
% nguon: Archive/SP_KhoLatex/ID6_old/1K3-9.tex #2660

    Phỏng vấn một số học sinh khối 11 về thời gian (giờ) ngủ của một buổi tối, thu được bảng số
    liệu ở bên.
    \begin{center}
        \begin{tabular}{|c|c|c|c|c|c|c|} 
            \hline
            Thời gian & Số học sinh nam & Số học sinh nữ  \\
            \hline
            [4;5) & $6$ & $4$\\
            \hline
            $[5;6)$ & $10$ & $8$\\
            \hline
            $[6; 7)$ & $13$ & $10$\\
            \hline
            $[7; 8)$ & $9$ & $11$\\
            \hline
            $[8; 9)$ & $7$ & $8$\\
            \hline
        \end{tabular}
    \end{center}
    \begin{enumerate}
        \item So sánh thời gian ngủ trung bình của các bạn học sinh nam và nữ.
        \item Hãy cho biết $75 \%$ học sinh khối $11$ ngủ ít nhất bao nhiêu giờ?
    \end{enumerate}
    \loigiai{
    \begin{enumerate}
        \item Ta có mẫu số liệu ghép nhóm
        \begin{center}
            \begin{tabular}{|c|c|c|c|c|c|c|} 
                \hline
                Thời gian & Giá trị đại diện & Số học sinh nam & Số học sinh nữ  \\
                \hline
                [4;5) & $4{,}5$ & $6$ & $4$\\
                \hline
                $[5;6)$ & $5{,}5$ & $10$ & $8$\\
                \hline
                $[6; 7)$ & $6{,}5$ & $13$ & $10$\\
                \hline
                $[7; 8)$ & $7{,}5$ & $9$ & $11$\\
                \hline
                $[8; 9)$ & $8{,}5$ & $7$ & $8$\\
                \hline
            \end{tabular}
        \end{center}
        Thời gian ngủ trung bình của các bạn nam là 
        $$\overline{x} = \dfrac{6 \cdot 4{,}5 + 10 \cdot 5{,}5 + 13 \cdot 6{,}5 + 9 \cdot 7{,}5 + 7 \cdot 8{,}5}{6 + 10 + 13 + 9 + 7} = 6{,}52.$$
        Thời gian ngủ trung bình của các bạn nam là 
        $$\overline{x} = \dfrac{4 \cdot 4{,}5 + 8 \cdot 5{,}5 + 10 \cdot 6{,}5 + 11 \cdot 7{,}5 + 8 \cdot 8{,}5}{4 + 8 + 10 + 11 + 8} = 6{,}77.$$
        Suy ra $6.77 > 6.52$. Như vậy thời gian ngủ trung bình của các bạn nữ nhiều hơn các bạn nam.
        \item Cỡ mẫu $n = 86$.\\
        Tứ phân vị thứ nhất $Q_1$ là $\dfrac{x_{21} + x_{22}}{2}$.\\
        Do $x_{21}$, $x_{22}$ đều thuộc nhóm $[5; 6)$ nên nhóm này chứa $Q_1$.\\
        Do đó $p = 2$, $a_2 = 5$; $m_2 = 18$; $m_1 = 10$; $a_3 - a_2 = 1$ và ta có 
        $$Q_1 = 5 + \dfrac{\dfrac{86}{4} - 10}{18} \cdot 1 = 5{,}64.$$
        Vậy $75 \%$ học sinh khối $11$ ngủ ít nhất $5{,}64$ giờ.
    \end{enumerate}
    }
\end{bt}

\begin{bt}%[1D5C2-3]
% nguon: Archive/SP_KhoLatex/ID6_old/1-GK1-KNTT-De15-NH23-24.tex #2713

	Số nguyện vọng đăng kí vào đại học của các bạn trong lớp được thống kê trong bảng sau
	\begin{center}
		\begin{tabular}{|c|c|c|c|c|}
		\hline
		Số nguyện vọng & $1-3$ & $4-6$ & $7-9$ & $10-12$\\
		\hline
		Số học sinh & $5$ & $18$ & $13$ & $7$\\
		\hline
	\end{tabular}
	\end{center}
\begin{itemize}
	\item[a)] Trung bình một bạn trong lớp dăng kí bao nhiêu nguyện vọng?
	\item[b)] Tìm các tứ phân vị của mẫu số liệu.
\end{itemize}
	\loigiai{
		\begin{itemize}
			\item[a)] Ta có cỡ mẫu là $n=5+18+13+7=43$.\\
			Do đó số trung bình của mẫu số liệu $x=\dfrac{5\cdot2+18\cdot5+13\cdot8+7\cdot11}{43}\approx6{,}53$.
			\item[b)] Hiệu chỉnh mẫu số liệu ta được bảng thống kê sau
		\begin{center}
				\begin{tabular}{|c|c|c|c|c|}
				\hline
				Số nguyện vọng& $\left[0{,}5;3{,}5\right)$ & $\left[3{,}5;6{,}5\right)$ & $\left[6{,}5;9{,}5\right)$ & $\left[9{,}5;12{,}5\right)$\\
				\hline
				Số học sinh& $5$ & $18$ & $13$ & $7$\\
				\hline
			\end{tabular}
		\end{center}
		\end{itemize}
	\begin{itemize}	
	\item[*] Vì $n=43$ nên trung vị là $x_{22}$ thuộc nhóm $\left[3{,}5;6{,}5\right)$, do đó $p=2$; $a_2=3{,}5$; $a_3=6{,}5$; $m_2=18$, $m_1=5$ và $a_3-a_2=3$. Do đó $Q_2=M_e=3{,}5+\dfrac{\dfrac{43}{2}-5}{18}\cdot3=6{,}25$.
	\item[*] Tứ phân vị thứ nhất $ Q_1 $ là $x_{11}$ thuộc nhóm $\left[3{,}5;6{,}5\right)$ do đó $p=2$; $a_2=3{,}5$; $a_3=6{,}5$; $m_2=18$; $m_1=5$ và $a_3-a_2=3$. Do đó $Q_1=3{,}5+\dfrac{\dfrac{43}{4}-5}{18}\cdot3=4{,}46$.
	\item[*] Tứ phân vị thứ ba $ Q_3 $ là $x_{33}$ thuộc nhóm $\left[6{,}5;9{,}5\right)$ do đó $p=3$; $a_3=6{,}5$; $a_4=9{,}5$; $m_1+m_2=23$ và $a_4-a_3=3$. Do đó $Q_3=6{,}5+\dfrac{\dfrac{3\cdot 43}{4}-23}{13}\cdot3=8{,}63$.
	\end{itemize}
	Vậy tứ phân vị của mẫu số số liệu là $Q_1=4{,}46$, $Q_2=6{,}25$ và $Q_3=8{,}63$.
	}
\end{bt}
